\section{待ち行列}

\subsection{平成23年度 問1-2}
\label{q:h23-reg-basic-1-2}

\noindent\textbf{原典：} sources/H23/basic.pdf、PDF物理1頁。

\textbf{公式正答：} 4（single）。

\par\textit{原典本文の校正と独自解説は未作成。分類のみの登録項目。}

\subsection{平成25年度 問1-5}
\label{q:h25-reg-basic-1-5}

\noindent\textbf{原典：} sources/H25/basic.pdf、PDF物理5頁。

\textbf{公式正答：} 4（single）。

\par\textit{原典本文の校正と独自解説は未作成。分類のみの登録項目。}

\subsection{平成27年度 問1-2}
\label{q:h27-reg-basic-1-2}

\noindent\textbf{原典：} sources/H27/basic.pdf、PDF物理2頁。

\textbf{公式正答：} 4（single）。

\par\textit{原典本文の校正と独自解説は未作成。分類のみの登録項目。}

\subsection{平成29年度 問1-1}
\label{q:h29-reg-basic-1-1}

\noindent\textbf{原典：} sources/H29/basic.pdf、PDF物理1頁。

\textbf{公式正答：} 5（single）。

\par\textit{原典本文の校正と独自解説は未作成。分類のみの登録項目。}

\subsection{令和元年度 問1-5}
\label{q:r01-reg-basic-1-5}

\noindent\textbf{原典：} sources/R01/basic.pdf、PDF物理5頁。

\textbf{公式正答：} 2（single）。

\par\textit{原典本文の校正と独自解説は未作成。分類のみの登録項目。}

\subsection{令和7年度 問1-5}
\label{q:r07-reg-basic-1-5}

\noindent\textbf{原典：} sources/R07/basic.pdf、PDF物理3頁。

\textbf{公式正答：} 5（single）。

\par\textit{原典本文の校正と独自解説は未作成。分類のみの登録項目。}

\section{工程管理・PERT}

\subsection{平成23年度 問1-3}
\label{q:h23-reg-basic-1-3}

\noindent\textbf{原典：} sources/H23/basic.pdf、PDF物理2頁。

\textbf{公式正答：} 3（single）。

\par\textit{原典本文の校正と独自解説は未作成。分類のみの登録項目。}

\subsection{平成25年度 問1-4}
\label{q:h25-reg-basic-1-4}

\noindent\textbf{原典：} sources/H25/basic.pdf、PDF物理4頁。

\textbf{公式正答：} 4（single）。

\par\textit{原典本文の校正と独自解説は未作成。分類のみの登録項目。}

\subsection{平成28年度 問1-4}
\label{q:h28-reg-basic-1-4}

\noindent\textbf{原典：} sources/H28/basic.pdf、PDF物理4頁。

\textbf{公式正答：} 4（single）。

\par\textit{原典本文の校正と独自解説は未作成。分類のみの登録項目。}

\subsection{平成30年度 問1-2}
\label{q:h30-reg-basic-1-2}

\noindent\textbf{原典：} sources/H30/basic.pdf、PDF物理2頁。

\textbf{公式正答：} 1（single）。

\par\textit{原典本文の校正と独自解説は未作成。分類のみの登録項目。}

\subsection{令和元年度再試験 問1-4}
\label{q:r01-ret-basic-1-4}

\noindent\textbf{原典：} sources/R01\_retest/basic.pdf、PDF物理4頁。

\textbf{公式正答：} 2（single）。

\par\textit{原典本文の校正と独自解説は未作成。分類のみの登録項目。}

\subsection{令和3年度 問1-3}
\label{q:r03-reg-basic-1-3}

\noindent\textbf{原典：} sources/R03/basic.pdf、PDF物理3頁。

\textbf{公式正答：} 2（single）。

\par\textit{原典本文の校正と独自解説は未作成。分類のみの登録項目。}

\subsection{令和7年度 問1-2}
\label{q:r07-reg-basic-1-2}

\noindent\textbf{原典：} sources/R07/basic.pdf、PDF物理2頁。

\textbf{公式正答：} 4（single）。

\par\textit{原典本文の校正と独自解説は未作成。分類のみの登録項目。}

\section{最適化・意思決定}

\subsection{平成23年度 問1-4}
\label{q:h23-reg-basic-1-4}

\noindent\textbf{原典：} sources/H23/basic.pdf、PDF物理3頁。

\textbf{公式正答：} 1（single）。

\par\textit{原典本文の校正と独自解説は未作成。分類のみの登録項目。}

\subsection{平成24年度 問1-5}
\label{q:h24-reg-basic-1-5}

\noindent\textbf{原典：} sources/H24/basic.pdf、PDF物理4頁。

\textbf{公式正答：} 5（single）。

\par\textit{原典本文の校正と独自解説は未作成。分類のみの登録項目。}

\subsection{平成25年度 問1-3}
\label{q:h25-reg-basic-1-3}

\noindent\textbf{原典：} sources/H25/basic.pdf、PDF物理3頁。

\textbf{公式正答：} 2（single）。

\par\textit{原典本文の校正と独自解説は未作成。分類のみの登録項目。}

\subsection{平成26年度 問1-2}
\label{q:h26-reg-basic-1-2}

\noindent\textbf{原典：} sources/H26/basic.pdf、PDF物理1頁。

\textbf{公式正答：} 4（single）。

\par\textit{原典本文の校正と独自解説は未作成。分類のみの登録項目。}

\subsection{平成28年度 問1-5}
\label{q:h28-reg-basic-1-5}

\noindent\textbf{原典：} sources/H28/basic.pdf、PDF物理5頁。

\textbf{公式正答：} 5（single）。

\par\textit{原典本文の校正と独自解説は未作成。分類のみの登録項目。}

\subsection{平成29年度 問1-3}
\label{q:h29-reg-basic-1-3}

\noindent\textbf{原典：} sources/H29/basic.pdf、PDF物理2頁。

\textbf{公式正答：} 2（single）。

\par\textit{原典本文の校正と独自解説は未作成。分類のみの登録項目。}

\subsection{平成30年度 問1-4}
\label{q:h30-reg-basic-1-4}

\noindent\textbf{原典：} sources/H30/basic.pdf、PDF物理4頁。

\textbf{公式正答：} 2（single）。

\par\textit{原典本文の校正と独自解説は未作成。分類のみの登録項目。}

\subsection{令和元年度 問1-1}
\label{q:r01-reg-basic-1-1}

\noindent\textbf{原典：} sources/R01/basic.pdf、PDF物理1頁。

\textbf{公式正答：} 5（single）。

\par\textit{原典本文の校正と独自解説は未作成。分類のみの登録項目。}

\subsection{令和2年度 問1-4}
\label{q:r02-reg-basic-1-4}

\noindent\textbf{原典：} sources/R02/basic.pdf、PDF物理4頁。

\textbf{公式正答：} 5（single）。

\par\textit{原典本文の校正と独自解説は未作成。分類のみの登録項目。}

\subsection{令和6年度 問1-5}
\label{q:r06-reg-basic-1-5}

\noindent\textbf{原典：} sources/R06/basic.pdf、PDF物理4頁。

\textbf{公式正答：} 4（single）。

\par\textit{原典本文の校正と独自解説は未作成。分類のみの登録項目。}

\section{品質管理・統計}

\subsection{平成24年度 問1-3}
\label{q:h24-reg-basic-1-3}

\noindent\textbf{原典：} sources/H24/basic.pdf、PDF物理3頁。

\textbf{公式正答：} 3（single）。

\par\textit{原典本文の校正と独自解説は未作成。分類のみの登録項目。}

\subsection{平成25年度 問1-6}
\label{q:h25-reg-basic-1-6}

\noindent\textbf{原典：} sources/H25/basic.pdf、PDF物理5頁。

\textbf{公式正答：} 3（single）。

\par\textit{原典本文の校正と独自解説は未作成。分類のみの登録項目。}

\subsection{平成26年度 問1-3}
\label{q:h26-reg-basic-1-3}

\noindent\textbf{原典：} sources/H26/basic.pdf、PDF物理2頁。

\textbf{公式正答：} 1（single）。

\par\textit{原典本文の校正と独自解説は未作成。分類のみの登録項目。}

\subsection{平成27年度 問1-5}
\label{q:h27-reg-basic-1-5}

\noindent\textbf{原典：} sources/H27/basic.pdf、PDF物理4頁。

\textbf{公式正答：} 1（single）。

\par\textit{原典本文の校正と独自解説は未作成。分類のみの登録項目。}

\subsection{平成27年度 問1-6}
\label{q:h27-reg-basic-1-6}

\noindent\textbf{原典：} sources/H27/basic.pdf、PDF物理4頁。

\textbf{公式正答：} 5（single）。

\par\textit{原典本文の校正と独自解説は未作成。分類のみの登録項目。}

\subsection{平成28年度 問1-2}
\label{q:h28-reg-basic-1-2}

\noindent\textbf{原典：} sources/H28/basic.pdf、PDF物理2頁。

\textbf{公式正答：} 4（single）。

\par\textit{原典本文の校正と独自解説は未作成。分類のみの登録項目。}

\subsection{平成30年度 問1-5}
\label{q:h30-reg-basic-1-5}

\noindent\textbf{原典：} sources/H30/basic.pdf、PDF物理5頁。

\textbf{公式正答：} 3（single）。

\par\textit{原典本文の校正と独自解説は未作成。分類のみの登録項目。}

\subsection{令和元年度 問1-3}
\label{q:r01-reg-basic-1-3}

\noindent\textbf{原典：} sources/R01/basic.pdf、PDF物理3頁。

\textbf{公式正答：} 5（single）。

\par\textit{原典本文の校正と独自解説は未作成。分類のみの登録項目。}

\subsection{令和2年度 問1-2}
\label{q:r02-reg-basic-1-2}

\noindent\textbf{原典：} sources/R02/basic.pdf、PDF物理2頁。

\textbf{公式正答：} 4（single）。

\par\textit{原典本文の校正と独自解説は未作成。分類のみの登録項目。}

\subsection{令和4年度 問1-3}
\label{q:r04-reg-basic-1-3}

\noindent\textbf{原典：} sources/R04/basic.pdf、PDF物理2頁。

\textbf{公式正答：} 3（single）。

\par\textit{原典本文の校正と独自解説は未作成。分類のみの登録項目。}

\subsection{令和5年度 問1-6}
\label{q:r05-reg-basic-1-6}

\noindent\textbf{原典：} sources/R05/basic.pdf、PDF物理4頁。

\textbf{公式正答：} 2（single）。

\par\textit{原典本文の校正と独自解説は未作成。分類のみの登録項目。}

\subsection{令和7年度 問1-6}
\label{q:r07-reg-basic-1-6}

\noindent\textbf{原典：} sources/R07/basic.pdf、PDF物理4頁。

\textbf{公式正答：} 3（single）。

\par\textit{原典本文の校正と独自解説は未作成。分類のみの登録項目。}

\section{論理式・命題}

\subsection{平成28年度 問1-6}
\label{q:h28-reg-basic-1-6}

\noindent\textbf{原典：} sources/H28/basic.pdf、PDF物理6頁。

\textbf{公式正答：} 1（single）。

\par\textit{原典本文の校正と独自解説は未作成。分類のみの登録項目。}

\section{数値表現}

\subsection{平成26年度 問1-6}
\label{q:h26-reg-basic-1-6}

\noindent\textbf{原典：} sources/H26/basic.pdf、PDF物理4頁。

\textbf{公式正答：} 3（single）。

\par\textit{原典本文の校正と独自解説は未作成。分類のみの登録項目。}

\subsection{平成29年度 問1-5}
\label{q:h29-reg-basic-1-5}

\noindent\textbf{原典：} sources/H29/basic.pdf、PDF物理4頁。

\textbf{公式正答：} 1（single）。

\par\textit{原典本文の校正と独自解説は未作成。分類のみの登録項目。}

\section{微積分・求積}

\subsection{令和元年度 問1-6}
\label{q:r01-reg-basic-1-6}

\noindent\textbf{原典：} sources/R01/basic.pdf、PDF物理6頁。

\textbf{公式正答：} 3（single）。

\par\textit{原典本文の校正と独自解説は未作成。分類のみの登録項目。}

\subsection{令和元年度再試験 問1-1}
\label{q:r01-ret-basic-1-1}

\noindent\textbf{原典：} sources/R01\_retest/basic.pdf、PDF物理1頁。

\textbf{公式正答：} 5（single）。

\par\textit{原典本文の校正と独自解説は未作成。分類のみの登録項目。}

\section{材料の応力・変形}

\subsection{令和2年度 問1-5}
\label{q:r02-reg-basic-1-5}

\noindent\textbf{原典：} sources/R02/basic.pdf、PDF物理5頁。

\textbf{公式正答：} 5（single）。

\par\textit{原典本文の校正と独自解説は未作成。分類のみの登録項目。}

\subsection{令和7年度 問1-4}
\label{q:r07-reg-basic-1-4}

\noindent\textbf{原典：} sources/R07/basic.pdf、PDF物理3頁。

\textbf{公式正答：} 1（single）。

\par\textit{原典本文の校正と独自解説は未作成。分類のみの登録項目。}

\section{環境法・循環・持続可能性}

\subsection{令和7年度 問1-1}
\label{q:r07-reg-basic-1-1}

\noindent\textbf{原典：} sources/R07/basic.pdf、PDF物理1頁。

\textbf{公式正答：} 3（single）。

\par\textit{原典本文の校正と独自解説は未作成。分類のみの登録項目。}

\section{修復系の定常稼働率}

\subsection{令和3年度 問1-4}
\label{q:r03-reg-basic-1-4}

\noindent\textbf{原典：} sources/R03/basic.pdf、PDF物理3頁。

\textbf{公式正答：} 3（single）。

\paragraph{問題の要旨} 平均運転時間と平均修復時間から、修復可能な装置の長期の稼働割合を導く。

\Needspace{4\baselineskip}\paragraph{確認した条件・記述（独自要約）}\mbox{}\par

\noindent\textbf{判断方向：} 全候補を比較し最も適切な語又は数値・順位を選ぶ。\par

\noindent\textbf{原典照合：} 25原典画像で22問の全候補、回路接続、FTAゲート、表の列順、数値・不等号を確認。\par

\noindent\textbf{前提：} 事象の独立性、非零根、修復モデル、法令時点と規準適用範囲は解法と疑義欄に区別する。\par

\paragraph{問われる核心} 平均運転時間と平均修復時間から、修復可能な装置の長期の稼働割合を導く。

\paragraph{必要知識} 平均運転時間、平均非稼働時間、交互修復の周期平均、信頼度R(t)と稼働率の相違。

\paragraph{考え方・正答を選ぶ理由（査読待ち初稿）} 試験で用いるMTBFを修復完了から次故障までの平均運転時間A、MTTRを平均修復時間Bとして読む。運転と修復を交互に繰り返す二状態系の長期割合は、有限平均と安定した周期モデルのもとで \(\mathcal A=E[T_{\mathrm{up}}]/(E[T_{\mathrm{up}}]+E[T_{\mathrm{down}}])=A/(A+B)\)、3である。指数分布の仮定はこの平均周期による割合には必須でない。待機・補給・計画停止を含める稼働率では分母の非稼働時間を拡張する。MTBFという名前だけで停止時間込みの暦時間間隔を代入しない。

\Needspace{4\baselineskip}\paragraph{選択肢の検討（査読待ち初稿）}\mbox{}\par

\noindent\textbf{1：不適合} \(A/(A-B)\)。周期は差ではなく運転と修復の和であり、B>0なら1を超える場合もある。\par

\noindent\textbf{2：不適合} \(B/(A-B)\)。周期の分母と、稼働時間の分子がともに違う。\par

\noindent\textbf{3：適合（記載条件内）} \(A/(A+B)\)。平均運転時間Aを、運転と修復の平均周期A+Bで割った稼働割合である。\par

\noindent\textbf{4：不適合} \(B/(A+B)\)。同じ二状態モデルでは非稼働率である。\par

\noindent\textbf{5：不適合} \(A/B\)。運転/修復の比は稼働の全時間割合ではなく1を超え得る。\par

\paragraph{再利用する考え方} 信頼度と稼働率、安全性と機能継続、規準に依存する係数を分ける。論理接続を事象に置き換え、積を使う前に独立性を確認する。

\paragraph{共通理論} \hyperref[theory:dependability:availability]{第\ref*{theory:dependability:availability}節}

\par 難易度 2（校準初稿）、優先度未評価。

\section{故障時の安全・縮退・機能継続}

\subsection{令和5年度 問1-5}
\label{q:r05-reg-basic-1-5}

\noindent\textbf{原典：} sources/R05/basic.pdf、PDF物理4頁。

\textbf{公式正答：} 1（single）。

\paragraph{問題の要旨} 故障時の安全確保、機能縮退、誤操作対策、故障下の機能継続を区別する。

\Needspace{4\baselineskip}\paragraph{確認した条件・記述（独自要約）}\mbox{}\par

\noindent\textbf{判断方向：} 全候補を比較し最も適切な語又は数値・順位を選ぶ。\par

\noindent\textbf{原典照合：} 25原典画像で22問の全候補、回路接続、FTAゲート、表の列順、数値・不等号を確認。\par

\noindent\textbf{前提：} 事象の独立性、非零根、修復モデル、法令時点と規準適用範囲は解法と疑義欄に区別する。\par

\paragraph{問われる核心} 故障時の安全確保、機能縮退、誤操作対策、故障下の機能継続を区別する。

\paragraph{必要知識} 機器故障、人為的誤操作、機能縮退、安全状態、耐故障能力と故障影響隠蔽の範囲。

\paragraph{考え方・正答を選ぶ理由（査読待ち初稿）} 故障しても安全側を保つ性質はフェールセーフである。故障又は故障の切迫時に影響機能を優先順位に従って徐々に終了するのがフェールソフトである。人為的な誤操作を重大な失敗へつなげない設計はフールプルーフ、想定された故障が存在しても機能を継続できる能力はフォールトトレランスである。四つをこの順に置く1を選ぶ。フォールトマスキングは故障影響を冗長性等で外へ現さない方式であり、故障の検出・隔離・回復を含む一般的な耐故障能力と同じ範囲ではない。安全に停止してよい系と停止自体が危険な系では、要求する安全状態を先に定める。

\Needspace{4\baselineskip}\paragraph{選択肢の検討（査読待ち初稿）}\mbox{}\par

\noindent\textbf{1：適合（記載条件内）} フェールセーフ・フェールソフト・フールプルーフ・フォールトトレランス。安全側、優先順位付き縮退、人為的誤り対策、故障下での機能継続の四記述に一致する。\par

\noindent\textbf{2：不適合} フェールセーフ・フェールソフト・フールプルーフ・フォールトマスキング。最後の一般的能力を、故障影響を外へ現さない一方式であるマスキングに限定している。\par

\noindent\textbf{3：不適合} フェールソフト・フォールトトレランス・フールプルーフ・フォールトマスキング。アは安全性を維持するフェールセーフ、イは段階的な機能終了のフェールソフトであり、エの一般的能力はトレランスである。\par

\noindent\textbf{4：不適合} フールプルーフ・フォールトトレランス・フェールソフト・フォールトマスキング。アの故障は人為的誤りのフールプルーフと異なり、イ・ウ・エの対応も違う。\par

\noindent\textbf{5：不適合} フールプルーフ・フェールセーフ・フェールソフト・フォールトトレランス。アはフェールセーフ、イはフェールソフト、ウはフールプルーフなので最初の三欄が違う。\par

\paragraph{再利用する考え方} 信頼度と稼働率、安全性と機能継続、規準に依存する係数を分ける。論理接続を事象に置き換え、積を使う前に独立性を確認する。

\paragraph{共通理論} \hyperref[theory:dependability:handling]{第\ref*{theory:dependability:handling}節}

\par 参考：\cite{design_mhlw_foolproof,design_mhlw_failsafe,design_nasa_tolerance}。

\par 難易度 2（校準初稿）、優先度未評価。

\section{故障の木と論理ゲート}

\subsection{令和元年度再試験 問1-3}
\label{q:r01-ret-basic-1-3}

\noindent\textbf{原典：} sources/R01\_retest/basic.pdf、PDF物理3頁。

\textbf{公式正答：} 1（single）。

\paragraph{問題の要旨} 独立な末端故障事象のFTAで、下位ゲートから最上位事象の発生確率を求める。

\Needspace{4\baselineskip}\paragraph{確認した条件・記述（独自要約）}\mbox{}\par

\noindent\textbf{判断方向：} 全候補を比較し最も適切な語又は数値・順位を選ぶ。\par

\noindent\textbf{原典照合：} 25原典画像で22問の全候補、回路接続、FTAゲート、表の列順、数値・不等号を確認。\par

\noindent\textbf{前提：} 事象の独立性、非零根、修復モデル、法令時点と規準適用範囲は解法と疑義欄に区別する。\par

\paragraph{問われる核心} 独立な末端故障事象のFTAで、下位ゲートから最上位事象の発生確率を求める。

\paragraph{必要知識} 直列は全要素成功の共通事象、並列は少なくとも一つ成功の和事象。独立時の積と補集合、確率範囲と丸め。

\paragraph{考え方・正答を選ぶ理由（査読待ち初稿）} 原典は各下流事象の独立を明記する。左枝の下位ORは \(1-0.9^2=0.19\)、その隣のANDは \(0.2^2=0.04\)、これらのORは \(1-(1-0.19)(1-0.04)=0.2224\)。右枝のANDは \(0.4^2=0.16\)。最上位はANDなので \(P=0.2224\cdot0.16=0.035584\)、近い0.036の1である。図は故障事象の木であり、ANDを直列系の「成功」と取り違えない。異なる枝に同じ末端事象を再使用する木では枝間の独立を別途点検する。

\Needspace{4\baselineskip}\paragraph{選択肢の検討（査読待ち初稿）}\mbox{}\par

\noindent\textbf{1：適合（記載条件内）} 0.036。導出値0.035584に最も近い。\par

\noindent\textbf{2：不適合} 0.038。導出値0.035584と一致せず、本問の接続・条件を満たす数値ではない。\par

\noindent\textbf{3：不適合} 0.233。導出値0.035584と一致せず、本問の接続・条件を満たす数値ではない。\par

\noindent\textbf{4：不適合} 0.641。導出値0.035584と一致せず、本問の接続・条件を満たす数値ではない。\par

\noindent\textbf{5：不適合} 0.804。導出値0.035584と一致せず、本問の接続・条件を満たす数値ではない。\par

\paragraph{再利用する考え方} 信頼度と稼働率、安全性と機能継続、規準に依存する係数を分ける。論理接続を事象に置き換え、積を使う前に独立性を確認する。

\paragraph{共通理論} \hyperref[theory:dependability:fta]{第\ref*{theory:dependability:fta}節}

\par 難易度 3（校準初稿）、優先度未評価。

\section{時間計画・状態監視・事後保全}

\subsection{平成26年度 問1-5}
\label{q:h26-reg-basic-1-5}

\noindent\textbf{原典：} sources/H26/basic.pdf、PDF物理3頁。

\textbf{公式正答：} 1（single）。

\paragraph{問題の要旨} 故障前後と保全時期の決定基準から、予防・事後・状態監視・時間計画の保全を対応させる。

\Needspace{4\baselineskip}\paragraph{確認した条件・記述（独自要約）}\mbox{}\par

\noindent\textbf{判断方向：} 全候補を比較し最も適切な語又は数値・順位を選ぶ。\par

\noindent\textbf{原典照合：} 25原典画像で22問の全候補、回路接続、FTAゲート、表の列順、数値・不等号を確認。\par

\noindent\textbf{前提：} 事象の独立性、非零根、修復モデル、法令時点と規準適用範囲は解法と疑義欄に区別する。\par

\paragraph{問われる核心} 故障前後と保全時期の決定基準から、予防・事後・状態監視・時間計画の保全を対応させる。

\paragraph{必要知識} 故障前/後、暦・運転時間の予定、状態測定による処置判断、定期点検と処置時期の区別。

\paragraph{考え方・正答を選ぶ理由（査読待ち初稿）} 故障を防ぐ目的で故障前に行う予防保全と、故障後に回復する事後保全を分ける。状態監視保全は点検値等から状態を評価して必要な処置を行う。時間計画保全では暦や運転時間等で予定を決め、定期的な保全や経時に基づく保全が対応する。本問の第三欄は状態監視、表の五欄は予防・事後・状態監視・定期・経時、1である。両年度は第三欄が異なるため、用語が似ているだけで候補を転用しない。状態を確認する点検が定期的であっても、それだけで状態に基づく処置が時間計画保全へ変わるわけではない。

\Needspace{4\baselineskip}\paragraph{選択肢の検討（査読待ち初稿）}\mbox{}\par

\noindent\textbf{1：適合（記載条件内）} 予防・事後・状態監視・定期・経時。故障前/後、保全時期の判断基準、定期と経時の対応が一致する。\par

\noindent\textbf{2：不適合} 定期・事後・経時・状態監視・予防。第1欄は予防であり定期ではない。第3欄は状態監視であり経時ではない。第4欄は定期であり状態監視ではない。第5欄は経時であり予防ではない。\par

\noindent\textbf{3：不適合} 状態監視・事後・予防・定期・経時。第1欄は予防であり状態監視ではない。第3欄は状態監視であり予防ではない。\par

\noindent\textbf{4：不適合} 定期・経時・予防・状態監視・事後。第1欄は予防であり定期ではない。第2欄は事後であり経時ではない。第3欄は状態監視であり予防ではない。第4欄は定期であり状態監視ではない。第5欄は経時であり事後ではない。\par

\noindent\textbf{5：不適合} 予防・定期・状態監視・経時・事後。第2欄は事後であり定期ではない。第4欄は定期であり経時ではない。第5欄は経時であり事後ではない。\par

\paragraph{再利用する考え方} 信頼度と稼働率、安全性と機能継続、規準に依存する係数を分ける。論理接続を事象に置き換え、積を使う前に独立性を確認する。

\paragraph{共通理論} \hyperref[theory:dependability:maintenance]{第\ref*{theory:dependability:maintenance}節}

\par 参考：\cite{maintenance_mlit_2013}。

\paragraph{関連問題} \hyperref[q:r01-ret-basic-1-6]{R01-RET-BASIC-1-6}

\par 難易度 2（校準初稿）、優先度未評価。

\subsection{令和元年度再試験 問1-6}
\label{q:r01-ret-basic-1-6}

\noindent\textbf{原典：} sources/R01\_retest/basic.pdf、PDF物理5頁。

\textbf{公式正答：} 5（single）。

\paragraph{問題の要旨} 故障前後と保全時期の決定基準から、予防・事後・状態監視・時間計画の保全を対応させる。

\Needspace{4\baselineskip}\paragraph{確認した条件・記述（独自要約）}\mbox{}\par

\noindent\textbf{判断方向：} 全候補を比較し最も適切な語又は数値・順位を選ぶ。\par

\noindent\textbf{原典照合：} 25原典画像で22問の全候補、回路接続、FTAゲート、表の列順、数値・不等号を確認。\par

\noindent\textbf{前提：} 事象の独立性、非零根、修復モデル、法令時点と規準適用範囲は解法と疑義欄に区別する。\par

\paragraph{問われる核心} 故障前後と保全時期の決定基準から、予防・事後・状態監視・時間計画の保全を対応させる。

\paragraph{必要知識} 故障前/後、暦・運転時間の予定、状態測定による処置判断、定期点検と処置時期の区別。

\paragraph{考え方・正答を選ぶ理由（査読待ち初稿）} 故障を防ぐ目的で故障前に行う予防保全と、故障後に回復する事後保全を分ける。状態監視保全は点検値等から状態を評価して必要な処置を行う。時間計画保全では暦や運転時間等で予定を決め、定期的な保全や経時に基づく保全が対応する。本問の第三欄は時間計画、表の五欄は予防・事後・時間計画・定期・経時、5である。両年度は第三欄が異なるため、用語が似ているだけで候補を転用しない。状態を確認する点検が定期的であっても、それだけで状態に基づく処置が時間計画保全へ変わるわけではない。

\Needspace{4\baselineskip}\paragraph{選択肢の検討（査読待ち初稿）}\mbox{}\par

\noindent\textbf{1：不適合} 予防・事後・劣化基準・状態監視・経時。第3欄は時間計画であり劣化基準ではない。第4欄は定期であり状態監視ではない。\par

\noindent\textbf{2：不適合} 状態監視・経時・時間計画・定期・予防。第1欄は予防であり状態監視ではない。第2欄は事後であり経時ではない。第5欄は経時であり予防ではない。\par

\noindent\textbf{3：不適合} 状態監視・事後・劣化基準・定期・経時。第1欄は予防であり状態監視ではない。第3欄は時間計画であり劣化基準ではない。\par

\noindent\textbf{4：不適合} 定期・経時・時間計画・状態監視・事後。第1欄は予防であり定期ではない。第2欄は事後であり経時ではない。第4欄は定期であり状態監視ではない。第5欄は経時であり事後ではない。\par

\noindent\textbf{5：適合（記載条件内）} 予防・事後・時間計画・定期・経時。故障前/後、保全時期の判断基準、定期と経時の対応が一致する。\par

\paragraph{再利用する考え方} 信頼度と稼働率、安全性と機能継続、規準に依存する係数を分ける。論理接続を事象に置き換え、積を使う前に独立性を確認する。

\paragraph{共通理論} \hyperref[theory:dependability:maintenance]{第\ref*{theory:dependability:maintenance}節}

\par 参考：\cite{maintenance_mlit_2013}。

\paragraph{関連問題} \hyperref[q:h26-reg-basic-1-5]{H26-REG-BASIC-1-5}

\par 難易度 2（校準初稿）、優先度未評価。

\section{単体法・双対と多目的最適化}

\subsection{令和元年度再試験 問1-2}
\label{q:r01-ret-basic-1-2}

\noindent\textbf{原典：} sources/R01\_retest/basic.pdf、PDF物理2頁。

\textbf{公式正答：} 3（single）。

\par\textbf{検討未確定：} 原文の「全ての目的関数を同時に改善できない」は全成分の厳密改善を意味すると弱パレートの条件であり通常の非支配性より弱い。初回/公式3と定義上の反例を別々に保持。

\paragraph{問題の要旨} 線形計画の解法と双対、多目的設計の相反関係とパレートの意味を整理する。

\Needspace{4\baselineskip}\paragraph{確認した条件・記述（独自要約）}\mbox{}\par

\noindent\textbf{判断方向：} 全候補を比較し最も適切な語又は数値・順位を選ぶ。\par

\noindent\textbf{原典照合：} 25原典画像で22問の全候補、回路接続、FTAゲート、表の列順、数値・不等号を確認。\par

\noindent\textbf{前提：} 事象の独立性、非零根、修復モデル、法令時点と規準適用範囲は解法と疑義欄に区別する。\par

\paragraph{問われる核心} 線形計画の解法と双対、多目的設計の相反関係とパレートの意味を整理する。

\paragraph{必要知識} 線形計画、基底交換、弱/強双対、可行集合、目的間の相反、支配と弱パレートの論理。

\paragraph{考え方・正答を選ぶ理由（査読待ち初稿）} 線形目的と線形制約の計画問題の代表解法はシンプレックス（単体）法であり、制約の重みを変数にする対応問題は双対である。複数の目的の改善が相反する関係はトレードオフで、他の可行解に支配されない解がパレート最適解である。表の意図はシンプレックス・双対・トレードオフ・パレートの3である。ただし最小化で「全てを同時に改善」を全成分の厳密減少と読むと弱パレートの条件になる。可行集合 \(\{(0,1),(0,0)\}\) では(0,1)は全成分の厳密減少を受けないが(0,0)に支配され、通常のパレート最適ではない。通常の定義は全成分で悪化せず少なくとも一成分を厳密改善する他解が存在しないことである。

\Needspace{4\baselineskip}\paragraph{選択肢の検討（査読待ち初稿）}\mbox{}\par

\noindent\textbf{1：不適合} シンプレックス・逆・トレードオン・パレート。線形計画で対応する問題は双対であり逆問題ではない。相反する改善はトレードオフである。\par

\noindent\textbf{2：不適合} シンプレックス・逆・トレードオフ・アクティブ。双対を逆とし、非支配解の呼称をアクティブとする点が違う。\par

\noindent\textbf{3：適合（記載条件内）} シンプレックス・双対・トレードオフ・パレート。単体法、双対、目的間の相反、意図されたパレートに対応する。ただし原文の全目的を同時改善できないという定義は弱パレートにも当てはまる。\par

\noindent\textbf{4：不適合} コンプレックス・逆・トレードオン・アクティブ。単体法のシンプレックス、双対、トレードオフ、パレートの四欄が異なる。\par

\noindent\textbf{5：不適合} コンプレックス・双対・トレードオン・パレート。単体法はシンプレックス、相反関係はトレードオフなので二欄が違う。\par

\paragraph{再利用する考え方} 信頼度と稼働率、安全性と機能継続、規準に依存する係数を分ける。論理接続を事象に置き換え、積を使う前に独立性を確認する。

\paragraph{共通理論} \hyperref[theory:dependability:optimization]{第\ref*{theory:dependability:optimization}節}

\par 参考：\cite{design_mit_lp}。

\par 難易度 3（校準初稿）、優先度未評価。

\section{製造物の欠陥と通常の安全性}

\subsection{平成25年度 問1-1}
\label{q:h25-reg-basic-1-1}

\noindent\textbf{原典：} sources/H25/basic.pdf、PDF物理1頁。

\textbf{公式正答：} 3（single）。

\par\textbf{検討未確定：} 原典の条文と平成30年消費者庁逐条解説の第1条・第2条を照合した。e-Gov現行本文は取得不能であり、出題年から現行までの全改正履歴は未照合。空欄の独立解と公式正答を保持し、現在の法全体が当時と同じだとはしない。

\paragraph{問題の要旨} 製造物責任法の保護目的と、通常有すべき安全性を欠くという欠陥の意味を整理する。

\Needspace{4\baselineskip}\paragraph{確認した条件・記述（独自要約）}\mbox{}\par

\noindent\textbf{判断方向：} 全候補を比較し最も適切な語又は数値・順位を選ぶ。\par

\noindent\textbf{原典照合：} 25原典画像で22問の全候補、回路接続、FTAゲート、表の列順、数値・不等号を確認。\par

\noindent\textbf{前提：} 事象の独立性、非零根、修復モデル、法令時点と規準適用範囲は解法と疑義欄に区別する。\par

\paragraph{問われる核心} 製造物責任法の保護目的と、通常有すべき安全性を欠くという欠陥の意味を整理する。

\paragraph{必要知識} 第1条の目的、第2条の製造物/欠陥、機能と安全性の違い、引渡時点と予見可能な使用。

\paragraph{考え方・正答を選ぶ理由（査読待ち初稿）} 第1条の対象は製造物、責任の根拠となる性状は欠陥、保護を図る相手は被害者である。第2条の欠陥は、製造物の特性、通常予見される使用形態、引渡時期等を考慮した通常有すべき安全性の不足である。従って製造物・欠陥・被害者・安全性の行を選ぶ。単なる故障、破損、性能や機能の不足が全てこの法の欠陥になるとは限らず、絶対に危険がないという要求でもない。被害者は購入者又は自然人だけへ限定しない。出題当時の原典条文と平成30年解説の比較に基づく空欄判断であり、時効等を含む法全体の現行解説とは区別する。

\Needspace{4\baselineskip}\paragraph{選択肢の検討（査読待ち初稿）}\mbox{}\par

\noindent\textbf{1：不適合} 製造物・欠陥・被害者・機能性。第4欄は安全性であり機能性ではない。\par

\noindent\textbf{2：不適合} 設計物・欠陥・製造者・安全性。第1欄は製造物であり設計物ではない。第3欄は被害者であり製造者ではない。\par

\noindent\textbf{3：適合（記載条件内）} 製造物・欠陥・被害者・安全性。第1条の対象・責任根拠・保護目的、第2条の欠陥の安全性概念に一致する。\par

\noindent\textbf{4：不適合} 設計物・破損・被害者・機能性。第1欄は製造物であり設計物ではない。第2欄は欠陥であり破損ではない。第4欄は安全性であり機能性ではない。\par

\noindent\textbf{5：不適合} 製造物・破損・製造者・機能性。第2欄は欠陥であり破損ではない。第3欄は被害者であり製造者ではない。第4欄は安全性であり機能性ではない。\par

\paragraph{再利用する考え方} 法定の欠陥と単なる不調を分け、設計・製造・指示警告の三側面から安全性を検討する。出題時点と現行解説の確認範囲を混同しない。

\paragraph{共通理論} \hyperref[theory:dependability:liability]{第\ref*{theory:dependability:liability}節}

\par 参考：\cite{design_caa_pl1,design_caa_pl2}。

\paragraph{関連問題} \hyperref[q:h27-reg-basic-1-4]{H27-REG-BASIC-1-4}

\par 難易度 2（校準初稿）、優先度未評価。

\subsection{平成27年度 問1-4}
\label{q:h27-reg-basic-1-4}

\noindent\textbf{原典：} sources/H27/basic.pdf、PDF物理3頁。

\textbf{公式正答：} 3（single）。

\par\textbf{検討未確定：} 原典の条文と平成30年消費者庁逐条解説の第1条・第2条を照合した。e-Gov現行本文は取得不能であり、出題年から現行までの全改正履歴は未照合。空欄の独立解と公式正答を保持し、現在の法全体が当時と同じだとはしない。

\paragraph{問題の要旨} 製造物責任法の保護目的と、通常有すべき安全性を欠くという欠陥の意味を整理する。

\Needspace{4\baselineskip}\paragraph{確認した条件・記述（独自要約）}\mbox{}\par

\noindent\textbf{判断方向：} 全候補を比較し最も適切な語又は数値・順位を選ぶ。\par

\noindent\textbf{原典照合：} 25原典画像で22問の全候補、回路接続、FTAゲート、表の列順、数値・不等号を確認。\par

\noindent\textbf{前提：} 事象の独立性、非零根、修復モデル、法令時点と規準適用範囲は解法と疑義欄に区別する。\par

\paragraph{問われる核心} 製造物責任法の保護目的と、通常有すべき安全性を欠くという欠陥の意味を整理する。

\paragraph{必要知識} 第1条の目的、第2条の製造物/欠陥、機能と安全性の違い、引渡時点と予見可能な使用。

\paragraph{考え方・正答を選ぶ理由（査読待ち初稿）} 第1条の対象は製造物、責任の根拠となる性状は欠陥、保護を図る相手は被害者である。第2条の欠陥は、製造物の特性、通常予見される使用形態、引渡時期等を考慮した通常有すべき安全性の不足である。従って製造物・欠陥・被害者・安全性の行を選ぶ。単なる故障、破損、性能や機能の不足が全てこの法の欠陥になるとは限らず、絶対に危険がないという要求でもない。被害者は購入者又は自然人だけへ限定しない。出題当時の原典条文と平成30年解説の比較に基づく空欄判断であり、時効等を含む法全体の現行解説とは区別する。

\Needspace{4\baselineskip}\paragraph{選択肢の検討（査読待ち初稿）}\mbox{}\par

\noindent\textbf{1：不適合} 製造物・欠陥・被害者・機能性。第4欄は安全性であり機能性ではない。\par

\noindent\textbf{2：不適合} 設計物・欠陥・製造者・安全性。第1欄は製造物であり設計物ではない。第3欄は被害者であり製造者ではない。\par

\noindent\textbf{3：適合（記載条件内）} 製造物・欠陥・被害者・安全性。第1条の対象・責任根拠・保護目的、第2条の欠陥の安全性概念に一致する。\par

\noindent\textbf{4：不適合} 設計物・破損・被害者・機能性。第1欄は製造物であり設計物ではない。第2欄は欠陥であり破損ではない。第4欄は安全性であり機能性ではない。\par

\noindent\textbf{5：不適合} 製造物・破損・製造者・機能性。第2欄は欠陥であり破損ではない。第3欄は被害者であり製造者ではない。第4欄は安全性であり機能性ではない。\par

\paragraph{再利用する考え方} 法定の欠陥と単なる不調を分け、設計・製造・指示警告の三側面から安全性を検討する。出題時点と現行解説の確認範囲を混同しない。

\paragraph{共通理論} \hyperref[theory:dependability:liability]{第\ref*{theory:dependability:liability}節}

\par 参考：\cite{design_caa_pl1,design_caa_pl2}。

\paragraph{関連問題} \hyperref[q:h25-reg-basic-1-1]{H25-REG-BASIC-1-1}, \hyperref[q:h30-reg-basic-1-6]{H30-REG-BASIC-1-6}

\par 難易度 2（校準初稿）、優先度未評価。

\subsection{平成30年度 問1-6}
\label{q:h30-reg-basic-1-6}

\noindent\textbf{原典：} sources/H30/basic.pdf、PDF物理5頁。

\textbf{公式正答：} 4（single）。

\par\textbf{検討未確定：} 原典の条文と平成30年消費者庁逐条解説の第1条・第2条を照合した。e-Gov現行本文は取得不能であり、出題年から現行までの全改正履歴は未照合。空欄の独立解と公式正答を保持し、現在の法全体が当時と同じだとはしない。

\paragraph{問題の要旨} 製造物責任法の保護目的と、通常有すべき安全性を欠くという欠陥の意味を整理する。

\Needspace{4\baselineskip}\paragraph{確認した条件・記述（独自要約）}\mbox{}\par

\noindent\textbf{判断方向：} 全候補を比較し最も適切な語又は数値・順位を選ぶ。\par

\noindent\textbf{原典照合：} 25原典画像で22問の全候補、回路接続、FTAゲート、表の列順、数値・不等号を確認。\par

\noindent\textbf{前提：} 事象の独立性、非零根、修復モデル、法令時点と規準適用範囲は解法と疑義欄に区別する。\par

\paragraph{問われる核心} 製造物責任法の保護目的と、通常有すべき安全性を欠くという欠陥の意味を整理する。

\paragraph{必要知識} 第1条の目的、第2条の製造物/欠陥、機能と安全性の違い、引渡時点と予見可能な使用。

\paragraph{考え方・正答を選ぶ理由（査読待ち初稿）} 第1条の対象は製造物、責任の根拠となる性状は欠陥、保護を図る相手は被害者である。第2条の欠陥は、製造物の特性、通常予見される使用形態、引渡時期等を考慮した通常有すべき安全性の不足である。従って製造物・欠陥・被害者・安全性の行を選ぶ。単なる故障、破損、性能や機能の不足が全てこの法の欠陥になるとは限らず、絶対に危険がないという要求でもない。被害者は購入者又は自然人だけへ限定しない。出題当時の原典条文と平成30年解説の比較に基づく空欄判断であり、時効等を含む法全体の現行解説とは区別する。

\Needspace{4\baselineskip}\paragraph{選択肢の検討（査読待ち初稿）}\mbox{}\par

\noindent\textbf{1：不適合} 製造物・故障・被害者・機能性。第2欄は欠陥であり故障ではない。第4欄は安全性であり機能性ではない。\par

\noindent\textbf{2：不適合} 設計物・欠陥・製造者・安全性。第1欄は製造物であり設計物ではない。第3欄は被害者であり製造者ではない。\par

\noindent\textbf{3：不適合} 設計物・破損・被害者・信頼性。第1欄は製造物であり設計物ではない。第2欄は欠陥であり破損ではない。第4欄は安全性であり信頼性ではない。\par

\noindent\textbf{4：適合（記載条件内）} 製造物・欠陥・被害者・安全性。第1条の対象・責任根拠・保護目的、第2条の欠陥の安全性概念に一致する。\par

\noindent\textbf{5：不適合} 製造物・破損・製造者・機能性。第2欄は欠陥であり破損ではない。第3欄は被害者であり製造者ではない。第4欄は安全性であり機能性ではない。\par

\paragraph{再利用する考え方} 法定の欠陥と単なる不調を分け、設計・製造・指示警告の三側面から安全性を検討する。出題時点と現行解説の確認範囲を混同しない。

\paragraph{共通理論} \hyperref[theory:dependability:liability]{第\ref*{theory:dependability:liability}節}

\par 参考：\cite{design_caa_pl1,design_caa_pl2}。

\paragraph{関連問題} \hyperref[q:h27-reg-basic-1-4]{H27-REG-BASIC-1-4}

\par 難易度 2（校準初稿）、優先度未評価。

\section{直列・並列の信頼度と独立性}

\subsection{平成23年度 問1-1}
\label{q:h23-reg-basic-1-1}

\noindent\textbf{原典：} sources/H23/basic.pdf、PDF物理1頁。

\textbf{公式正答：} 2（single）。

\par\textbf{検討未確定：} 原典は要素/路線の独立を明記しない。周辺信頼度だけでは結合確率が決まらず、積の式は独立モデルを補った条件付き解である。公式正答と独立初回を保持する。

\paragraph{問題の要旨} 既存の二路線を直列経路とし、直結路線の追加で全体の信頼度を指定値へ上げる。

\Needspace{4\baselineskip}\paragraph{確認した条件・記述（独自要約）}\mbox{}\par

\noindent\textbf{判断方向：} 全候補を比較し最も適切な語又は数値・順位を選ぶ。\par

\noindent\textbf{原典照合：} 25原典画像で22問の全候補、回路接続、FTAゲート、表の列順、数値・不等号を確認。\par

\noindent\textbf{前提：} 事象の独立性、非零根、修復モデル、法令時点と規準適用範囲は解法と疑義欄に区別する。\par

\paragraph{問われる核心} 既存の二路線を直列経路とし、直結路線の追加で全体の信頼度を指定値へ上げる。

\paragraph{必要知識} 直列は全要素成功の共通事象、並列は少なくとも一つ成功の和事象。独立時の積と補集合、確率範囲と丸め。

\paragraph{考え方・正答を選ぶ理由（査読待ち初稿）} 東京から大阪まで名古屋経由の成功事象は二路線の両成功である。独立なら \(R_0=0.8\cdot0.6=0.48\)。直結の信頼度を \(x\) とし、旧経路との独立も仮定すると \(R=1-(1-0.48)(1-x)\)。\(R=0.8\) より \(0.52(1-x)=0.2\)、\(x=1-0.2/0.52=0.6153846\)、最も近い0.62の2である。路線の共通災害で同時不通になる場合にこの独立式をそのまま使わない。

\Needspace{4\baselineskip}\paragraph{選択肢の検討（査読待ち初稿）}\mbox{}\par

\noindent\textbf{1：不適合} 0.52。導出値0.6153846と一致せず、本問の接続・条件を満たす数値ではない。\par

\noindent\textbf{2：適合（記載条件内）} 0.62。導出値0.6153846に最も近い。\par

\noindent\textbf{3：不適合} 0.72。導出値0.6153846と一致せず、本問の接続・条件を満たす数値ではない。\par

\noindent\textbf{4：不適合} 0.82。導出値0.6153846と一致せず、本問の接続・条件を満たす数値ではない。\par

\noindent\textbf{5：不適合} 0.92。導出値0.6153846と一致せず、本問の接続・条件を満たす数値ではない。\par

\paragraph{再利用する考え方} 信頼度と稼働率、安全性と機能継続、規準に依存する係数を分ける。論理接続を事象に置き換え、積を使う前に独立性を確認する。

\paragraph{共通理論} \hyperref[theory:dependability:series]{第\ref*{theory:dependability:series}節}, \hyperref[theory:dependability:parallel]{第\ref*{theory:dependability:parallel}節}

\par 参考：\cite{reliability_nist_series,reliability_nist_parallel}。

\par 難易度 2（校準初稿）、優先度未評価。

\subsection{平成24年度 問1-1}
\label{q:h24-reg-basic-1-1}

\noindent\textbf{原典：} sources/H24/basic.pdf、PDF物理1頁。

\textbf{公式正答：} 4（single）。

\par\textbf{検討未確定：} 原典は要素/路線の独立を明記しない。周辺信頼度だけでは結合確率が決まらず、積の式は独立モデルを補った条件付き解である。公式正答と独立初回を保持する。

\paragraph{問題の要旨} 並列二要素と直列一要素からなる経路へ、もう一経路を並列接続した信頼度を求める。

\Needspace{4\baselineskip}\paragraph{確認した条件・記述（独自要約）}\mbox{}\par

\noindent\textbf{判断方向：} 全候補を比較し最も適切な語又は数値・順位を選ぶ。\par

\noindent\textbf{原典照合：} 25原典画像で22問の全候補、回路接続、FTAゲート、表の列順、数値・不等号を確認。\par

\noindent\textbf{前提：} 事象の独立性、非零根、修復モデル、法令時点と規準適用範囲は解法と疑義欄に区別する。\par

\paragraph{問われる核心} 並列二要素と直列一要素からなる経路へ、もう一経路を並列接続した信頼度を求める。

\paragraph{必要知識} 直列は全要素成功の共通事象、並列は少なくとも一つ成功の和事象。独立時の積と補集合、確率範囲と丸め。

\paragraph{考え方・正答を選ぶ理由（査読待ち初稿）} 図を内側から読む。独立なら \(R_{12}=1-(1-0.9)(1-0.8)=0.98\)、これと \(R_3=0.5\) の直列は \(0.49\)。\(R_4=0.7\) がこの経路と並列なので \(R=1-(1-0.49)(1-0.7)=0.8474\)。近い0.85の4である。最外側の並列を、全四要素の直列や全四要素の並列と読み替えない。

\Needspace{4\baselineskip}\paragraph{選択肢の検討（査読待ち初稿）}\mbox{}\par

\noindent\textbf{1：不適合} 0.25。導出値0.8474と一致せず、本問の接続・条件を満たす数値ではない。\par

\noindent\textbf{2：不適合} 0.60。導出値0.8474と一致せず、本問の接続・条件を満たす数値ではない。\par

\noindent\textbf{3：不適合} 0.73。導出値0.8474と一致せず、本問の接続・条件を満たす数値ではない。\par

\noindent\textbf{4：適合（記載条件内）} 0.85。導出値0.8474に最も近い。\par

\noindent\textbf{5：不適合} 0.91。導出値0.8474と一致せず、本問の接続・条件を満たす数値ではない。\par

\paragraph{再利用する考え方} 信頼度と稼働率、安全性と機能継続、規準に依存する係数を分ける。論理接続を事象に置き換え、積を使う前に独立性を確認する。

\paragraph{共通理論} \hyperref[theory:dependability:series]{第\ref*{theory:dependability:series}節}, \hyperref[theory:dependability:parallel]{第\ref*{theory:dependability:parallel}節}

\par 参考：\cite{reliability_nist_series,reliability_nist_parallel}。

\paragraph{関連問題} \hyperref[q:h26-reg-basic-1-4]{H26-REG-BASIC-1-4}

\par 難易度 2（校準初稿）、優先度未評価。

\subsection{平成26年度 問1-4}
\label{q:h26-reg-basic-1-4}

\noindent\textbf{原典：} sources/H26/basic.pdf、PDF物理3頁。

\textbf{公式正答：} 3（single）。

\par\textbf{検討未確定：} 原典は要素/路線の独立を明記しない。周辺信頼度だけでは結合確率が決まらず、積の式は独立モデルを補った条件付き解である。公式正答と独立初回を保持する。

\paragraph{問題の要旨} 四要素の直並列系と同じ信頼度を持つ三要素直列系の、単体信頼度を逆算する。

\Needspace{4\baselineskip}\paragraph{確認した条件・記述（独自要約）}\mbox{}\par

\noindent\textbf{判断方向：} 全候補を比較し最も適切な語又は数値・順位を選ぶ。\par

\noindent\textbf{原典照合：} 25原典画像で22問の全候補、回路接続、FTAゲート、表の列順、数値・不等号を確認。\par

\noindent\textbf{前提：} 事象の独立性、非零根、修復モデル、法令時点と規準適用範囲は解法と疑義欄に区別する。\par

\paragraph{問われる核心} 四要素の直並列系と同じ信頼度を持つ三要素直列系の、単体信頼度を逆算する。

\paragraph{必要知識} 直列は全要素成功の共通事象、並列は少なくとも一つ成功の和事象。独立時の積と補集合、確率範囲と丸め。

\paragraph{考え方・正答を選ぶ理由（査読待ち初稿）} 独立ならA系の両端の0.9は直列、中央の二つの0.9は並列である。\(R_A=0.9^2[1-(1-0.9)^2]=0.8019\)。B系は同じ信頼度 \(x\) の三要素の直列なので \(x^3=0.8019\)、\(x=0.8019^{1/3}=0.9290521\)。近い0.93の3である。初回記録の概算0.928はそのまま保存し、ここで精算値を示す。

\Needspace{4\baselineskip}\paragraph{選択肢の検討（査読待ち初稿）}\mbox{}\par

\noindent\textbf{1：不適合} 0.87。導出値0.9290521と一致せず、本問の接続・条件を満たす数値ではない。\par

\noindent\textbf{2：不適合} 0.90。導出値0.9290521と一致せず、本問の接続・条件を満たす数値ではない。\par

\noindent\textbf{3：適合（記載条件内）} 0.93。導出値0.9290521に最も近い。\par

\noindent\textbf{4：不適合} 0.96。導出値0.9290521と一致せず、本問の接続・条件を満たす数値ではない。\par

\noindent\textbf{5：不適合} 0.99。導出値0.9290521と一致せず、本問の接続・条件を満たす数値ではない。\par

\paragraph{再利用する考え方} 信頼度と稼働率、安全性と機能継続、規準に依存する係数を分ける。論理接続を事象に置き換え、積を使う前に独立性を確認する。

\paragraph{共通理論} \hyperref[theory:dependability:series]{第\ref*{theory:dependability:series}節}, \hyperref[theory:dependability:parallel]{第\ref*{theory:dependability:parallel}節}

\par 参考：\cite{reliability_nist_series,reliability_nist_parallel}。

\paragraph{関連問題} \hyperref[q:h24-reg-basic-1-1]{H24-REG-BASIC-1-1}

\par 難易度 2（校準初稿）、優先度未評価。

\subsection{平成27年度 問1-1}
\label{q:h27-reg-basic-1-1}

\noindent\textbf{原典：} sources/H27/basic.pdf、PDF物理1頁。

\textbf{公式正答：} 5（single）。

\paragraph{問題の要旨} 同じ単体信頼度を持つ単独・二要素直列・二要素並列の大小を比較する。

\Needspace{4\baselineskip}\paragraph{確認した条件・記述（独自要約）}\mbox{}\par

\noindent\textbf{判断方向：} 全候補を比較し最も適切な語又は数値・順位を選ぶ。\par

\noindent\textbf{原典照合：} 25原典画像で22問の全候補、回路接続、FTAゲート、表の列順、数値・不等号を確認。\par

\noindent\textbf{前提：} 事象の独立性、非零根、修復モデル、法令時点と規準適用範囲は解法と疑義欄に区別する。\par

\paragraph{問われる核心} 同じ単体信頼度を持つ単独・二要素直列・二要素並列の大小を比較する。

\paragraph{必要知識} 独立二事象、共通/和事象、0<r<1での不等式。

\paragraph{考え方・正答を選ぶ理由（査読待ち初稿）} 独立性は原典に明記される。単独 \(A=0.9\)、二要素直列 \(B=0.9^2=0.81\)、二要素並列 \(C=1-(1-0.9)^2=0.99\) なので \(B<A<C\)、5である。一般に \(0<r<1\) なら \(r^2<r<2r-r^2\) であり、両端r=\allowbreak{}0,1では等しくなる。

\Needspace{4\baselineskip}\paragraph{選択肢の検討（査読待ち初稿）}\mbox{}\par

\noindent\textbf{1：不適合} B>A>C。直列Bが単独Aより高いという最初の関係が逆である。\par

\noindent\textbf{2：不適合} B>A=\allowbreak{}C。BはAより小さく、AとCも等しくない。\par

\noindent\textbf{3：不適合} B=\allowbreak{}A=\allowbreak{}C。三者は0.81、0.90、0.99で等しくない。\par

\noindent\textbf{4：不適合} B<A=\allowbreak{}C。B<Aは合うがA=\allowbreak{}Cは違う。\par

\noindent\textbf{5：適合（記載条件内）} B<A<C。0.81<0.90<0.99に一致する。\par

\paragraph{再利用する考え方} 信頼度と稼働率、安全性と機能継続、規準に依存する係数を分ける。論理接続を事象に置き換え、積を使う前に独立性を確認する。

\paragraph{共通理論} \hyperref[theory:dependability:series]{第\ref*{theory:dependability:series}節}, \hyperref[theory:dependability:parallel]{第\ref*{theory:dependability:parallel}節}

\par 参考：\cite{reliability_nist_series,reliability_nist_parallel}。

\paragraph{関連問題} \hyperref[q:r03-reg-basic-1-2]{R03-REG-BASIC-1-2}

\par 難易度 2（校準初稿）、優先度未評価。

\subsection{平成28年度 問1-1}
\label{q:h28-reg-basic-1-1}

\noindent\textbf{原典：} sources/H28/basic.pdf、PDF物理1頁。

\textbf{公式正答：} 4（single）。

\par\textbf{検討未確定：} 原典は要素/路線の独立を明記しない。周辺信頼度だけでは結合確率が決まらず、積の式は独立モデルを補った条件付き解である。公式正答と独立初回を保持する。さらに非零信頼度の条件がなくx=\allowbreak{}0も解である。公式4はx>0を補った解として示す。

\paragraph{問題の要旨} 共通信頼度の未知要素を含む二つの直並列系が等しいとき、全体の信頼度を求める。

\Needspace{4\baselineskip}\paragraph{確認した条件・記述（独自要約）}\mbox{}\par

\noindent\textbf{判断方向：} 全候補を比較し最も適切な語又は数値・順位を選ぶ。\par

\noindent\textbf{原典照合：} 25原典画像で22問の全候補、回路接続、FTAゲート、表の列順、数値・不等号を確認。\par

\noindent\textbf{前提：} 事象の独立性、非零根、修復モデル、法令時点と規準適用範囲は解法と疑義欄に区別する。\par

\paragraph{問われる核心} 共通信頼度の未知要素を含む二つの直並列系が等しいとき、全体の信頼度を求める。

\paragraph{必要知識} 直列は全要素成功の共通事象、並列は少なくとも一つ成功の和事象。独立時の積と補集合、確率範囲と丸め。

\paragraph{考え方・正答を選ぶ理由（査読待ち初稿）} 独立なら \(R_A=0.9^2[1-(1-x)^2]=0.81(2x-x^2)\)、\(R_B=0.95x^2\)。等式は \(1.62x-1.76x^2=0\)。非零の根を選ぶ場合 \(x=1.62/1.76=0.9204545\)、全体系 \(R=0.95x^2=0.8048747\) なので0.805の4である。ただし原典はx>0を明示せず、x=\allowbreak{}0でも等式を満たしR=\allowbreak{}0となる。未知単体xと問われた全体系Rも分ける。

\Needspace{4\baselineskip}\paragraph{選択肢の検討（査読待ち初稿）}\mbox{}\par

\noindent\textbf{1：不適合} 0.857。導出値0.8048747と一致せず、本問の接続・条件を満たす数値ではない。\par

\noindent\textbf{2：不適合} 0.839。導出値0.8048747と一致せず、本問の接続・条件を満たす数値ではない。\par

\noindent\textbf{3：不適合} 0.822。導出値0.8048747と一致せず、本問の接続・条件を満たす数値ではない。\par

\noindent\textbf{4：適合（記載条件内）} 0.805。導出値0.8048747に最も近い。\par

\noindent\textbf{5：不適合} 0.787。導出値0.8048747と一致せず、本問の接続・条件を満たす数値ではない。\par

\paragraph{再利用する考え方} 信頼度と稼働率、安全性と機能継続、規準に依存する係数を分ける。論理接続を事象に置き換え、積を使う前に独立性を確認する。

\paragraph{共通理論} \hyperref[theory:dependability:series]{第\ref*{theory:dependability:series}節}, \hyperref[theory:dependability:parallel]{第\ref*{theory:dependability:parallel}節}

\par 参考：\cite{reliability_nist_series,reliability_nist_parallel}。

\paragraph{関連問題} \hyperref[q:h30-reg-basic-1-1]{H30-REG-BASIC-1-1}

\par 難易度 3（校準初稿）、優先度未評価。

\subsection{平成30年度 問1-1}
\label{q:h30-reg-basic-1-1}

\noindent\textbf{原典：} sources/H30/basic.pdf、PDF物理1頁。

\textbf{公式正答：} 3（single）。

\par\textbf{検討未確定：} 原典は要素/路線の独立を明記しない。周辺信頼度だけでは結合確率が決まらず、積の式は独立モデルを補った条件付き解である。公式正答と独立初回を保持する。

\paragraph{問題の要旨} 未知単体信頼度を含む直並列系と三要素直列系の等式から、系の信頼度を求める。

\Needspace{4\baselineskip}\paragraph{確認した条件・記述（独自要約）}\mbox{}\par

\noindent\textbf{判断方向：} 全候補を比較し最も適切な語又は数値・順位を選ぶ。\par

\noindent\textbf{原典照合：} 25原典画像で22問の全候補、回路接続、FTAゲート、表の列順、数値・不等号を確認。\par

\noindent\textbf{前提：} 事象の独立性、非零根、修復モデル、法令時点と規準適用範囲は解法と疑義欄に区別する。\par

\paragraph{問われる核心} 未知単体信頼度を含む直並列系と三要素直列系の等式から、系の信頼度を求める。

\paragraph{必要知識} 直列は全要素成功の共通事象、並列は少なくとも一つ成功の和事象。独立時の積と補集合、確率範囲と丸め。

\paragraph{考え方・正答を選ぶ理由（査読待ち初稿）} 独立なら \(R_A=0.95^2(2x-x^2)=0.9025(2x-x^2)\)、\(R_B=x^3\)。原典の \(x>0\) を用いて割ると \(x^2+0.9025x-1.805=0\)。正根は \(x=(-0.9025+\sqrt{0.9025^2+4\cdot1.805})/2=0.9660102\)、もう一根は負で不適である。問われる系の信頼度は \(x^3=0.9014573\)、0.901の3である。0.966は単体の値であり全体値ではない。

\Needspace{4\baselineskip}\paragraph{選択肢の検討（査読待ち初稿）}\mbox{}\par

\noindent\textbf{1：不適合} 0.835。導出値0.9014573と一致せず、本問の接続・条件を満たす数値ではない。\par

\noindent\textbf{2：不適合} 0.857。導出値0.9014573と一致せず、本問の接続・条件を満たす数値ではない。\par

\noindent\textbf{3：適合（記載条件内）} 0.901。導出値0.9014573に最も近い。\par

\noindent\textbf{4：不適合} 0.945。導出値0.9014573と一致せず、本問の接続・条件を満たす数値ではない。\par

\noindent\textbf{5：不適合} 0.966。導出値0.9014573と一致せず、本問の接続・条件を満たす数値ではない。\par

\paragraph{再利用する考え方} 信頼度と稼働率、安全性と機能継続、規準に依存する係数を分ける。論理接続を事象に置き換え、積を使う前に独立性を確認する。

\paragraph{共通理論} \hyperref[theory:dependability:series]{第\ref*{theory:dependability:series}節}, \hyperref[theory:dependability:parallel]{第\ref*{theory:dependability:parallel}節}

\par 参考：\cite{reliability_nist_series,reliability_nist_parallel}。

\paragraph{関連問題} \hyperref[q:h28-reg-basic-1-1]{H28-REG-BASIC-1-1}

\par 難易度 3（校準初稿）、優先度未評価。

\subsection{令和2年度 問1-6}
\label{q:r02-reg-basic-1-6}

\noindent\textbf{原典：} sources/R02/basic.pdf、PDF物理6頁。

\textbf{公式正答：} 3（single）。

\paragraph{問題の要旨} 同じ部品の並列数を増やし、直列部品を含む系で要求信頼度を初めて満たす個数を求める。

\Needspace{4\baselineskip}\paragraph{確認した条件・記述（独自要約）}\mbox{}\par

\noindent\textbf{判断方向：} 全候補を比較し最も適切な語又は数値・順位を選ぶ。\par

\noindent\textbf{原典照合：} 25原典画像で22問の全候補、回路接続、FTAゲート、表の列順、数値・不等号を確認。\par

\noindent\textbf{前提：} 事象の独立性、非零根、修復モデル、法令時点と規準適用範囲は解法と疑義欄に区別する。\par

\paragraph{問われる核心} 同じ部品の並列数を増やし、直列部品を含む系で要求信頼度を初めて満たす個数を求める。

\paragraph{必要知識} 直列は全要素成功の共通事象、並列は少なくとも一つ成功の和事象。独立時の積と補集合、確率範囲と丸め。

\paragraph{考え方・正答を選ぶ理由（査読待ち初稿）} 独立性が明記される。信頼度0.7のn個並列と0.95の一要素が直列であるから \(R_n=0.95(1-0.3^n)\)。\(R_n\ge0.94\) は \(0.3^n\le1/95\)。\(n\ge2\) の整数を比較すると \(R_2=0.8645\)、\(R_3=0.92435\)、\(R_4=0.942305\)、\(R_5=0.9476915\)。単調増加なので最小は4、候補3である。上限は0.95であり0.94は達成可能である。

\Needspace{4\baselineskip}\paragraph{選択肢の検討（査読待ち初稿）}\mbox{}\par

\noindent\textbf{1：不適合} 2。導出値最小n=\allowbreak{}4と一致せず、本問の接続・条件を満たす数値ではない。\par

\noindent\textbf{2：不適合} 3。導出値最小n=\allowbreak{}4と一致せず、本問の接続・条件を満たす数値ではない。\par

\noindent\textbf{3：適合（記載条件内）} 4。導出値最小n=\allowbreak{}4に最も近い。\par

\noindent\textbf{4：不適合} 5。導出値最小n=\allowbreak{}4と一致せず、本問の接続・条件を満たす数値ではない。\par

\noindent\textbf{5：不適合} nをいくつにしても信頼度0.94以上にならない。導出値最小n=\allowbreak{}4と一致せず、本問の接続・条件を満たす数値ではない。\par

\paragraph{再利用する考え方} 信頼度と稼働率、安全性と機能継続、規準に依存する係数を分ける。論理接続を事象に置き換え、積を使う前に独立性を確認する。

\paragraph{共通理論} \hyperref[theory:dependability:series]{第\ref*{theory:dependability:series}節}, \hyperref[theory:dependability:parallel]{第\ref*{theory:dependability:parallel}節}

\par 参考：\cite{reliability_nist_series,reliability_nist_parallel}。

\paragraph{関連問題} \hyperref[q:r05-reg-basic-1-4]{R05-REG-BASIC-1-4}

\par 難易度 3（校準初稿）、優先度未評価。

\subsection{令和3年度 問1-2}
\label{q:r03-reg-basic-1-2}

\noindent\textbf{原典：} sources/R03/basic.pdf、PDF物理2頁。

\textbf{公式正答：} 2（single）。

\paragraph{問題の要旨} 同じ三要素の接続順序を変えた五つの系の信頼度を、成功経路から比較する。

\Needspace{4\baselineskip}\paragraph{確認した条件・記述（独自要約）}\mbox{}\par

\noindent\textbf{判断方向：} 全候補を比較し最も適切な語又は数値・順位を選ぶ。\par

\noindent\textbf{原典照合：} 25原典画像で22問の全候補、回路接続、FTAゲート、表の列順、数値・不等号を確認。\par

\noindent\textbf{前提：} 事象の独立性、非零根、修復モデル、法令時点と規準適用範囲は解法と疑義欄に区別する。\par

\paragraph{問われる核心} 同じ三要素の接続順序を変えた五つの系の信頼度を、成功経路から比較する。

\paragraph{必要知識} 五つの接続の識別、二要素直列経路と一要素並列経路、単体と系の値の区別。

\paragraph{考え方・正答を選ぶ理由（査読待ち初稿）} 独立を前提とする原典の接続を式へ写す。\(A=1-(1-0.9\cdot0.8)(1-0.7)=0.916\)、\(B=1-(1-0.9\cdot0.7)(1-0.8)=0.926\)、\(C=1-(1-0.8\cdot0.7)(1-0.9)=0.956\)。\(D=0.8[1-(1-0.9)(1-0.7)]=0.776\)、\(E=0.9[1-(1-0.8)(1-0.7)]=0.846\)。従ってC>B>A>E>D、2である。単体の値の順だけを見ず、どの要素が一つの独立経路として全体の停止を防げるかを読む。

\Needspace{4\baselineskip}\paragraph{選択肢の検討（査読待ち初稿）}\mbox{}\par

\noindent\textbf{1：不適合} C>B>E>A>D。A=\allowbreak{}0.916はE=\allowbreak{}0.846より高く、その順序が逆である。\par

\noindent\textbf{2：適合（記載条件内）} C>B>A>E>D。C=\allowbreak{}0.956、B=\allowbreak{}0.926、A=\allowbreak{}0.916、E=\allowbreak{}0.846、D=\allowbreak{}0.776の順位に一致する。\par

\noindent\textbf{3：不適合} C>E>B>D>A。EはBより低く、DはAより低い。\par

\noindent\textbf{4：不適合} E>D>A>B>C。EとDは前三系より低いので上位に置けない。\par

\noindent\textbf{5：不適合} E>D>C>B>A。EとDをCより上位に置く点が違う。\par

\paragraph{再利用する考え方} 信頼度と稼働率、安全性と機能継続、規準に依存する係数を分ける。論理接続を事象に置き換え、積を使う前に独立性を確認する。

\paragraph{共通理論} \hyperref[theory:dependability:series]{第\ref*{theory:dependability:series}節}, \hyperref[theory:dependability:parallel]{第\ref*{theory:dependability:parallel}節}

\par 参考：\cite{reliability_nist_series,reliability_nist_parallel}。

\paragraph{関連問題} \hyperref[q:h27-reg-basic-1-1]{H27-REG-BASIC-1-1}

\par 難易度 3（校準初稿）、優先度未評価。

\section{多数決冗長系の信頼度}

\subsection{令和5年度 問1-4}
\label{q:r05-reg-basic-1-4}

\noindent\textbf{原典：} sources/R05/basic.pdf、PDF物理3頁。

\textbf{公式正答：} 4（single）。

\paragraph{問題の要旨} 三要素のうち二つ以上が動けば成功する多数決系の信頼度の区間を求める。

\Needspace{4\baselineskip}\paragraph{確認した条件・記述（独自要約）}\mbox{}\par

\noindent\textbf{判断方向：} 全候補を比較し最も適切な語又は数値・順位を選ぶ。\par

\noindent\textbf{原典照合：} 25原典画像で22問の全候補、回路接続、FTAゲート、表の列順、数値・不等号を確認。\par

\noindent\textbf{前提：} 事象の独立性、非零根、修復モデル、法令時点と規準適用範囲は解法と疑義欄に区別する。\par

\paragraph{問われる核心} 三要素のうち二つ以上が動けば成功する多数決系の信頼度の区間を求める。

\paragraph{必要知識} 直列は全要素成功の共通事象、並列は少なくとも一つ成功の和事象。独立時の積と補集合、確率範囲と丸め。

\paragraph{考え方・正答を選ぶ理由（査読待ち初稿）} 独立な三要素の単体信頼度は0.7である。成功事象は「ちょうど二つ成功」と「三つ成功」の排反和であるから \(R={3\choose2}0.7^2(1-0.7)+0.7^3=0.441+0.343=0.784\)。\(0.75\le R<0.80\) の4である。少なくとも一つ成功する全並列、全三つ必要な直列、単に二要素だけを選ぶ直列とは異なる。判定器自身の故障は所与モデルに含まれておらず、その信頼性を実装で無条件に保証しない。

\Needspace{4\baselineskip}\paragraph{選択肢の検討（査読待ち初稿）}\mbox{}\par

\noindent\textbf{1：不適合} \(0.90\le R\le1.00\)。導出値0.784と一致せず、本問の接続・条件を満たす数値ではない。\par

\noindent\textbf{2：不適合} \(0.85\le R<0.90\)。導出値0.784と一致せず、本問の接続・条件を満たす数値ではない。\par

\noindent\textbf{3：不適合} \(0.80\le R<0.85\)。導出値0.784と一致せず、本問の接続・条件を満たす数値ではない。\par

\noindent\textbf{4：適合（記載条件内）} \(0.75\le R<0.80\)。導出値0.784に最も近い。\par

\noindent\textbf{5：不適合} \(0.70\le R<0.75\)。導出値0.784と一致せず、本問の接続・条件を満たす数値ではない。\par

\paragraph{再利用する考え方} 信頼度と稼働率、安全性と機能継続、規準に依存する係数を分ける。論理接続を事象に置き換え、積を使う前に独立性を確認する。

\paragraph{共通理論} \hyperref[theory:dependability:voting]{第\ref*{theory:dependability:voting}節}

\par 参考：\cite{reliability_nist_voting}。

\paragraph{関連問題} \hyperref[q:r02-reg-basic-1-6]{R02-REG-BASIC-1-6}

\par 難易度 2（校準初稿）、優先度未評価。

\section{用途・基準時点ごとの安全係数}

\subsection{平成23年度 問1-5}
\label{q:h23-reg-basic-1-5}

\noindent\textbf{原典：} sources/H23/basic.pdf、PDF物理4頁。

\textbf{公式正答：} 4（single）。

\par\textbf{検討未確定：} 航空宇宙と薬品の対象が広すぎ、四値の全称命題は検証できない。主索4.0は平成25年改正反映表で確認したが2011当時版の全条件未照合。公式/独立4を保持して適用範囲を限定。

\paragraph{問題の要旨} 異なる用途で使われる安全率の代表値を比較し、それぞれの定義と条件を区別する。

\Needspace{4\baselineskip}\paragraph{確認した条件・記述（独自要約）}\mbox{}\par

\noindent\textbf{判断方向：} 全候補を比較し最も適切な語又は数値・順位を選ぶ。\par

\noindent\textbf{原典照合：} 25原典画像で22問の全候補、回路接続、FTAゲート、表の列順、数値・不等号を確認。\par

\noindent\textbf{前提：} 事象の独立性、非零根、修復モデル、法令時点と規準適用範囲は解法と疑義欄に区別する。\par

\paragraph{問われる核心} 異なる用途で使われる安全率の代表値を比較し、それぞれの定義と条件を区別する。

\paragraph{必要知識} 限界/終局荷重比、切断荷重/最大荷重、無毒性量/許容摂取量、主索と端部、設置/使用時の区別。

\paragraph{考え方・正答を選ぶ理由（査読待ち初稿）} 代表値の意図を、航空機の構造1.5、玉掛けワイヤロープ6、摂取量の毒性評価100、使用時のエレベーター主索4として読むと4である。2011年米国25.303は航空機の限界荷重に原則1.5を掛ける規定であり、ロケット全体の普遍規準ではない。クレーン等安全規則213条は玉掛けワイヤロープの切断荷重/最大荷重を6以上とする。食品安全委員会の2008年説明は残留農薬等のADIを無毒性量から通常100で割る扱いであって治療薬全般の定数ではない。国交省の平成25年改正反映告示表では、限定された主索の常時使用時4.0、設置時5.0、端部は別値である。分母・対象・時点が異なる数値を同じ物理量の絶対的安全余裕として比較しない。

\Needspace{4\baselineskip}\paragraph{選択肢の検討（査読待ち初稿）}\mbox{}\par

\noindent\textbf{1：不適合} 100・6・1000・4。航空機の代表1.5、摂取量評価の代表100という対応と異なる。\par

\noindent\textbf{2：不適合} 1.5・4・100・6。玉掛けワイヤロープは6以上であり4ではない。使用時主索の表値も4.0であって6ではない。\par

\noindent\textbf{3：不適合} 100・4・1000・6。航空機、玉掛け、摂取量評価、使用時主索の四代表値との対応が違う。\par

\noindent\textbf{4：適合（記載条件内）} 1.5・6・100・4。意図された四代表値1.5・6・100・4に対応する。ただし航空宇宙/薬品という広い対象を一律の規定値とする根拠にはしない。\par

\noindent\textbf{5：不適合} 10・4・100・6。航空機の代表は1.5、玉掛けは6、対象告示表の使用時主索は4.0なので三欄が違う。\par

\paragraph{再利用する考え方} 信頼度と稼働率、安全性と機能継続、規準に依存する係数を分ける。論理接続を事象に置き換え、積を使う前に独立性を確認する。

\paragraph{共通理論} \hyperref[theory:dependability:factor]{第\ref*{theory:dependability:factor}節}

\par 参考：\cite{design_faa_2011,design_mhlw_crane,design_fsc_adi,design_mlit_elevator}。

\paragraph{関連問題} \hyperref[q:h29-reg-basic-1-2]{H29-REG-BASIC-1-2}

\par 難易度 3（校準初稿）、優先度未評価。

\subsection{平成27年度 問1-3}
\label{q:h27-reg-basic-1-3}

\noindent\textbf{原典：} sources/H27/basic.pdf、PDF物理2頁。

\textbf{公式正答：} 5（single）。

\par\textbf{検討未確定：} 安全率の一般的な2以上という説明は基準強度・材料・荷重・規準で変わり普遍の下限ではない。公式/独立5を保持して代表的な許容応力モデルと区別する。

\paragraph{問題の要旨} 許容応力の定義、安全率の代表値、線形弾性の法則と棒の伸びを対応させる。

\Needspace{4\baselineskip}\paragraph{確認した条件・記述（独自要約）}\mbox{}\par

\noindent\textbf{判断方向：} 全候補を比較し最も適切な語又は数値・順位を選ぶ。\par

\noindent\textbf{原典照合：} 25原典画像で22問の全候補、回路接続、FTAゲート、表の列順、数値・不等号を確認。\par

\noindent\textbf{前提：} 事象の独立性、非零根、修復モデル、法令時点と規準適用範囲は解法と疑義欄に区別する。\par

\paragraph{問われる核心} 許容応力の定義、安全率の代表値、線形弾性の法則と棒の伸びを対応させる。

\paragraph{必要知識} 基準強さ/安全率、フックの法則、一定断面棒の伸び、固定する長さ・弾性係数・断面積。

\paragraph{考え方・正答を選ぶ理由（査読待ち初稿）} 許容応力を基準の極限強さから定めるモデルは \(\sigma_{\mathrm{allow}}=\sigma_{\mathrm{u}}/n\)。線形弾性域の \(\sigma=E\varepsilon\) はフックの法則であり、一定断面・均一材料の軸荷重棒で \(\Delta L=FL/(EA)\)。L,E,Aを固定すると伸びは荷重Fに比例する。表の÷・2以上・フック・荷重の5である。ただし設問の一般的なnが2以上という代表的説明を、航空機やすべての設計規準へ適用する下限値にしない。

\Needspace{4\baselineskip}\paragraph{選択肢の検討（査読待ち初稿）}\mbox{}\par

\noindent\textbf{1：不適合} ÷・2以上・エネルギー・断面積。比例の法則はフックであり、同じL,E,Aの伸びは荷重に比例して断面積には反比例する。\par

\noindent\textbf{2：不適合} ÷・2以上・フック・断面積。最後の比例先は荷重であり断面積ではない。\par

\noindent\textbf{3：不適合} ×・1以下・フック・荷重。基準強度に安全率を掛けるのではなく割る。1以下は強度を十分下回る許容応力を作らない。\par

\noindent\textbf{4：不適合} ×・1以下・エネルギー・断面積。演算、代表値、法則、比例先がいずれも想定された対応と異なる。\par

\noindent\textbf{5：適合（記載条件内）} ÷・2以上・フック・荷重。基準強度を割る、設問の代表2以上、フックの法則、荷重比例の四欄が一致する。\par

\paragraph{再利用する考え方} 信頼度と稼働率、安全性と機能継続、規準に依存する係数を分ける。論理接続を事象に置き換え、積を使う前に独立性を確認する。

\paragraph{共通理論} \hyperref[theory:dependability:factor]{第\ref*{theory:dependability:factor}節}

\par 参考：\cite{struct_mit_beam,design_faa_2011}。

\par 難易度 2（校準初稿）、優先度未評価。

\subsection{平成29年度 問1-2}
\label{q:h29-reg-basic-1-2}

\noindent\textbf{原典：} sources/H29/basic.pdf、PDF物理1頁。

\textbf{公式正答：} 5（single）。

\par\textbf{検討未確定：} アの航空機/ロケット、ウの摂取薬品の範囲が規準を特定しない。100>6>1.5は限定された代表例の順序として保持し普遍命題にしない。

\paragraph{問題の要旨} 航空宇宙・玉掛け・薬品の安全率の代表値の順序を、用途の違いを含めて比較する。

\Needspace{4\baselineskip}\paragraph{確認した条件・記述（独自要約）}\mbox{}\par

\noindent\textbf{判断方向：} 全候補を比較し最も適切な語又は数値・順位を選ぶ。\par

\noindent\textbf{原典照合：} 25原典画像で22問の全候補、回路接続、FTAゲート、表の列順、数値・不等号を確認。\par

\noindent\textbf{前提：} 事象の独立性、非零根、修復モデル、法令時点と規準適用範囲は解法と疑義欄に区別する。\par

\paragraph{問われる核心} 航空宇宙・玉掛け・薬品の安全率の代表値の順序を、用途の違いを含めて比較する。

\paragraph{必要知識} 用途別の安全係数定義、基準荷重と毒性評価、代表値と規準値の差。

\paragraph{考え方・正答を選ぶ理由（査読待ち初稿）} アを航空機構造の原則1.5、イを玉掛けワイヤロープの6、ウをADI算定などの通常100という代表例として読むと、ウ>イ>ア、5である。数値の大きさは同じ破壊荷重比を使った横断的な安全順位ではない。毒性評価では種差と個体差等を補う不確実係数であり、航空機では限界荷重と終局荷重、玉掛けでは切断荷重と最大荷重の比である。治療薬、ロケットの全てに一律の代表値を適用しない。

\Needspace{4\baselineskip}\paragraph{選択肢の検討（査読待ち初稿）}\mbox{}\par

\noindent\textbf{1：不適合} ア>イ>ウ。意図された代表値ではウが最も大きいので逆である。\par

\noindent\textbf{2：不適合} イ>ウ>ア。100のウは6のイより大きい。\par

\noindent\textbf{3：不適合} ウ>ア>イ。1.5のアは6のイより小さい。\par

\noindent\textbf{4：不適合} ア>ウ>イ。1.5のアを100のウより大きくする点が違う。\par

\noindent\textbf{5：適合（記載条件内）} ウ>イ>ア。限定した代表例100>6>1.5に一致する。\par

\paragraph{再利用する考え方} 信頼度と稼働率、安全性と機能継続、規準に依存する係数を分ける。論理接続を事象に置き換え、積を使う前に独立性を確認する。

\paragraph{共通理論} \hyperref[theory:dependability:factor]{第\ref*{theory:dependability:factor}節}

\par 参考：\cite{design_faa_2011,design_mhlw_crane,design_fsc_adi}。

\paragraph{関連問題} \hyperref[q:h23-reg-basic-1-5]{H23-REG-BASIC-1-5}

\par 難易度 2（校準初稿）、優先度未評価。

\subsection{令和7年度 問1-3}
\label{q:r07-reg-basic-1-3}

\noindent\textbf{原典：} sources/R07/basic.pdf、PDF物理2頁。

\textbf{公式正答：} 2（single）。

\par\textbf{検討未確定：} 2以上と1.3～2.0は設問の一般的代表値である。全材料・全用途に成立する統一規準を一次資料で確認していない。公式/独立2を保持し普遍の安全条件としない。

\paragraph{問題の要旨} 許容応力を基準強さから定める演算と、静的強度・疲労限度に対する代表安全率を区別する。

\Needspace{4\baselineskip}\paragraph{確認した条件・記述（独自要約）}\mbox{}\par

\noindent\textbf{判断方向：} 全候補を比較し最も適切な語又は数値・順位を選ぶ。\par

\noindent\textbf{原典照合：} 25原典画像で22問の全候補、回路接続、FTAゲート、表の列順、数値・不等号を確認。\par

\noindent\textbf{前提：} 事象の独立性、非零根、修復モデル、法令時点と規準適用範囲は解法と疑義欄に区別する。\par

\paragraph{問われる核心} 許容応力を基準強さから定める演算と、静的強度・疲労限度に対する代表安全率を区別する。

\paragraph{必要知識} 許容応力の次元、基準強さの選定、引張強さと疲労限度、代表安全率と規準。

\paragraph{考え方・正答を選ぶ理由（査読待ち初稿）} 原典の分数は基準強さを安全率nで割る \(\sigma_{\mathrm{allow}}=\sigma_{\mathrm{ref}}/n\) である。nは無次元なので結果は応力の次元を持つ。引張強さを基準とする一般的代表値2以上、疲労限度を基準とする代表範囲1.3～2.0を表に置く2である。13～20という十倍の範囲へ小数点を落とさない。一方、実設計では降伏・座屈・疲労寿命・使用条件と指定規準を確認する必要があり、この代表範囲を全用途の法定下限とはしない。

\Needspace{4\baselineskip}\paragraph{選択肢の検討（査読待ち初稿）}\mbox{}\par

\noindent\textbf{1：不適合} 基準強さ/n・1以下・13～20。Bの1以下と疲労側13～20が設問の代表値に一致しない。\par

\noindent\textbf{2：適合（記載条件内）} 基準強さ/n・2以上・1.3～2.0。許容応力の次元と演算、設問で想定する静的2以上・疲労1.3～2.0に対応する。\par

\noindent\textbf{3：不適合} 基準強さ/n・2以上・13～20。Cは小数点を落とした13～20であり1.3～2.0ではない。\par

\noindent\textbf{4：不適合} n/基準強さ・1以下・1.3～2.0。Aは応力の逆数となり許容応力の次元でなく、Bも違う。\par

\noindent\textbf{5：不適合} n/基準強さ・2以上・13～20。Aの次元とCの小数点が違う。\par

\paragraph{再利用する考え方} 信頼度と稼働率、安全性と機能継続、規準に依存する係数を分ける。論理接続を事象に置き換え、積を使う前に独立性を確認する。

\paragraph{共通理論} \hyperref[theory:dependability:factor]{第\ref*{theory:dependability:factor}節}

\par 参考：\cite{struct_mit_beam,struct_nims_data,design_faa_2011}。

\par 難易度 2（校準初稿）、優先度未評価。

\section{断面二次モーメント・断面係数と寸法依存}

\subsection{令和4年度 問1-5}
\label{q:r04-reg-basic-1-5}

\noindent\textbf{原典：} sources/R04/basic.pdf、PDF物理3頁。

\textbf{公式正答：} 3（single）。

\par\textbf{検討未確定：} 径/高さを変える比較でq,L,E（矩形は幅b）を固定する前提が必要。自重のみでqが断面積に比例して変わる円梁なら応力はdの逆1乗、たわみは逆2乗となる。試験の通常の固定荷重解と区別。

\paragraph{問題の要旨} 等分布荷重を受ける単純支持梁の断面寸法と、最大曲げ応力・たわみのべき依存性を導く。

\Needspace{4\baselineskip}\paragraph{確認した条件・記述（独自要約）}\mbox{}\par

\noindent\textbf{判断方向：} 空欄の適切な語・数値又は正誤の組合せ。全候補の順序は各肢に記録。\par

\noindent\textbf{原典照合：} 16原典画像で17問の全候補、図、数値、不等号、表の列順を確認。\par

\noindent\textbf{前提：} 解法に記録した弾性範囲、荷重条件、幾何、独立性を明示する。\par

\paragraph{問われる核心} 等分布荷重を受ける単純支持梁の断面寸法と、最大曲げ応力・たわみのべき依存性を導く。

\paragraph{必要知識} 固定するq,L,E、円断面IとZ、矩形の幅固定、曲げモーメント、梁の線形小たわみ。

\paragraph{考え方・正答を選ぶ理由（査読待ち初稿）} 同じq,L,Eの円断面梁では \(M_{\max}=qL^2/8\)、\(I=\pi d^4/64\)、\(Z=I/(d/2)=\pi d^3/32\)。従って \(\sigma_{\max}=M_{\max}/Z\propto d^{-3}\)、\(\delta_{\max}=5qL^4/(384EI)\propto d^{-4}\)。矩形断面は幅bを固定して \(Z=bh^2/6\) より \(\sigma_{\max}\propto h^{-2}\)。A3、B反比例、C4、D2の3となる。

\Needspace{4\baselineskip}\paragraph{選択肢の検討（査読待ち初稿）}\mbox{}\par

\noindent\textbf{1：不適合} 3・比例・4・3。列2は反比例であり比例ではない。列4は2であり3ではない。\par

\noindent\textbf{2：不適合} 4・比例・3・2。列1は3であり4ではない。列2は反比例であり比例ではない。列3は4であり3ではない。\par

\noindent\textbf{3：適合（前提内）} 3・反比例・4・2。曲げ応力は断面係数Zの逆数、変位は断面二次モーメントIの逆数という二つの依存性に適合する。\par

\noindent\textbf{4：不適合} 4・反比例・3・3。列1は3であり4ではない。列3は4であり3ではない。列4は2であり3ではない。\par

\noindent\textbf{5：不適合} 3・反比例・4・3。列4は2であり3ではない。\par

\paragraph{再利用する考え方} 材料の強さ、構造の安定、変形による機能低下を分ける。寸法だけを変えるとき何を固定したかを明示し、同じ理論でも候補番号を年度間で流用しない。

\paragraph{共通理論} \hyperref[theory:structural:beam]{第\ref*{theory:structural:beam}節}, \hyperref[theory:structural:scaling]{第\ref*{theory:structural:scaling}節}

\par 参考：\cite{struct_mit_beam,struct_mit_buckling,struct_mlit_limit}。

\paragraph{関連問題} \hyperref[q:r06-reg-basic-1-3]{R06-REG-BASIC-1-3}

\par 難易度 3（校準初稿）、優先度未評価。

\subsection{令和6年度 問1-3}
\label{q:r06-reg-basic-1-3}

\noindent\textbf{原典：} sources/R06/basic.pdf、PDF物理2頁。

\textbf{公式正答：} 1（single）。

\par\textbf{検討未確定：} 径/高さを変える比較でq,L,E（矩形は幅b）を固定する前提が必要。自重のみでqが断面積に比例して変わる円梁なら応力はdの逆1乗、たわみは逆2乗となる。試験の通常の固定荷重解と区別。

\paragraph{問題の要旨} 等分布荷重を受ける単純支持梁の断面寸法と、最大曲げ応力・たわみのべき依存性を導く。

\Needspace{4\baselineskip}\paragraph{確認した条件・記述（独自要約）}\mbox{}\par

\noindent\textbf{判断方向：} 空欄の適切な語・数値又は正誤の組合せ。全候補の順序は各肢に記録。\par

\noindent\textbf{原典照合：} 16原典画像で17問の全候補、図、数値、不等号、表の列順を確認。\par

\noindent\textbf{前提：} 解法に記録した弾性範囲、荷重条件、幾何、独立性を明示する。\par

\paragraph{問われる核心} 等分布荷重を受ける単純支持梁の断面寸法と、最大曲げ応力・たわみのべき依存性を導く。

\paragraph{必要知識} 固定するq,L,E、円断面IとZ、矩形の幅固定、曲げモーメント、梁の線形小たわみ。

\paragraph{考え方・正答を選ぶ理由（査読待ち初稿）} 同じq,L,Eの円断面単純支持梁で \(I=\pi d^4/64\)、\(Z=\pi d^3/32\) である。最大モーメントは \(qL^2/8\)、最大たわみは \(5qL^4/(384EI)\)。従って \(\sigma_{\max}\propto d^{-3}\)、\(\delta_{\max}\propto d^{-4}\) であり、3乗に反比例・4乗に反比例の1を選ぶ。

\Needspace{4\baselineskip}\paragraph{選択肢の検討（査読待ち初稿）}\mbox{}\par

\noindent\textbf{1：適合（前提内）} 3・反比例・4・反比例。曲げ応力は断面係数Zの逆数、変位は断面二次モーメントIの逆数という二つの依存性に適合する。\par

\noindent\textbf{2：不適合} 4・比例・4・反比例。列1は3であり4ではない。列2は反比例であり比例ではない。\par

\noindent\textbf{3：不適合} 4・反比例・3・反比例。列1は3であり4ではない。列3は4であり3ではない。\par

\noindent\textbf{4：不適合} 4・比例・3・比例。列1は3であり4ではない。列2は反比例であり比例ではない。列3は4であり3ではない。列4は反比例であり比例ではない。\par

\noindent\textbf{5：不適合} 3・反比例・4・比例。列4は反比例であり比例ではない。\par

\paragraph{再利用する考え方} 材料の強さ、構造の安定、変形による機能低下を分ける。寸法だけを変えるとき何を固定したかを明示し、同じ理論でも候補番号を年度間で流用しない。

\paragraph{共通理論} \hyperref[theory:structural:beam]{第\ref*{theory:structural:beam}節}, \hyperref[theory:structural:scaling]{第\ref*{theory:structural:scaling}節}

\par 参考：\cite{struct_mit_beam,struct_mit_buckling,struct_mlit_limit}。

\paragraph{関連問題} \hyperref[q:r04-reg-basic-1-5]{R04-REG-BASIC-1-5}

\par 難易度 3（校準初稿）、優先度未評価。

\section{圧縮安定とオイラー座屈}

\subsection{平成25年度 問1-2}
\label{q:h25-reg-basic-1-2}

\noindent\textbf{原典：} sources/H25/basic.pdf、PDF物理2頁。

\textbf{公式正答：} 5（single）。

\paragraph{問題の要旨} 薄い定規の引張と圧縮の差を座屈で説明し、許容応力と安全率を対応させる。

\Needspace{4\baselineskip}\paragraph{確認した条件・記述（独自要約）}\mbox{}\par

\noindent\textbf{判断方向：} 空欄の適切な語・数値又は正誤の組合せ。全候補の順序は各肢に記録。\par

\noindent\textbf{原典照合：} 16原典画像で17問の全候補、図、数値、不等号、表の列順を確認。\par

\noindent\textbf{前提：} 解法に記録した弾性範囲、荷重条件、幾何、独立性を明示する。\par

\paragraph{問われる核心} 薄い定規の引張と圧縮の差を座屈で説明し、許容応力と安全率を対応させる。

\paragraph{必要知識} 材料強度と座屈荷重の相違、圧縮の安定性、許容応力＝強度/安全率。

\paragraph{考え方・正答を選ぶ理由（査読待ち初稿）} 引張では直線形状が安定しやすいが、圧縮では小さな横たわみが増幅され座屈する。Aは引張、Bは圧縮、Cは座屈、Dは安全率である。強度を \(\sigma_{\mathrm{str}}\)、安全率を \(n>1\) とすると \(\sigma_{\mathrm{allow}}=\sigma_{\mathrm{str}}/n\)。圧縮の座屈荷重はこの材料強度だけで決まらない。表の並びに従い本問の候補番号を選ぶ。

\Needspace{4\baselineskip}\paragraph{選択肢の検討（査読待ち初稿）}\mbox{}\par

\noindent\textbf{1：不適合} 引張・圧縮・剥離・安全率。横へたわむ圧縮不安定は座屈であり、層がはがれる剥離ではない。\par

\noindent\textbf{2：不適合} 圧縮・引張・剥離・安全率。薄い定規は引張より圧縮の安定性を失いやすく、AとBが逆である。横へたわむ圧縮不安定は座屈であり、層がはがれる剥離ではない。\par

\noindent\textbf{3：不適合} 圧縮・引張・剥離・ヤング率。薄い定規は引張より圧縮の安定性を失いやすく、AとBが逆である。横へたわむ圧縮不安定は座屈であり、層がはがれる剥離ではない。強度を割る無次元比は安全率であり、応力単位のヤング率ではない。\par

\noindent\textbf{4：不適合} 圧縮・引張・座屈・ヤング率。薄い定規は引張より圧縮の安定性を失いやすく、AとBが逆である。強度を割る無次元比は安全率であり、応力単位のヤング率ではない。\par

\noindent\textbf{5：適合（前提内）} 引張・圧縮・座屈・安全率。引張・圧縮・座屈・安全率が対応する。形の安定と材料強度は別の照査である。\par

\paragraph{再利用する考え方} 材料の強さ、構造の安定、変形による機能低下を分ける。寸法だけを変えるとき何を固定したかを明示し、同じ理論でも候補番号を年度間で流用しない。

\paragraph{共通理論} \hyperref[theory:structural:buckling]{第\ref*{theory:structural:buckling}節}, \hyperref[theory:structural:allowable]{第\ref*{theory:structural:allowable}節}

\par 参考：\cite{struct_mit_beam,struct_mit_buckling,struct_mlit_limit}。

\paragraph{関連問題} \hyperref[q:h28-reg-basic-1-3]{H28-REG-BASIC-1-3}, \hyperref[q:r03-reg-basic-1-5]{R03-REG-BASIC-1-5}

\par 難易度 2（校準初稿）、優先度未評価。

\subsection{平成28年度 問1-3}
\label{q:h28-reg-basic-1-3}

\noindent\textbf{原典：} sources/H28/basic.pdf、PDF物理3頁。

\textbf{公式正答：} 4（single）。

\paragraph{問題の要旨} 薄い定規の引張と圧縮の差を座屈で説明し、許容応力と安全率を対応させる。

\Needspace{4\baselineskip}\paragraph{確認した条件・記述（独自要約）}\mbox{}\par

\noindent\textbf{判断方向：} 空欄の適切な語・数値又は正誤の組合せ。全候補の順序は各肢に記録。\par

\noindent\textbf{原典照合：} 16原典画像で17問の全候補、図、数値、不等号、表の列順を確認。\par

\noindent\textbf{前提：} 解法に記録した弾性範囲、荷重条件、幾何、独立性を明示する。\par

\paragraph{問われる核心} 薄い定規の引張と圧縮の差を座屈で説明し、許容応力と安全率を対応させる。

\paragraph{必要知識} 材料強度と座屈荷重の相違、圧縮の安定性、許容応力＝強度/安全率。

\paragraph{考え方・正答を選ぶ理由（査読待ち初稿）} 引張では直線形状が安定しやすいが、圧縮では小さな横たわみが増幅され座屈する。Aは引張、Bは圧縮、Cは座屈、Dは安全率である。強度を \(\sigma_{\mathrm{str}}\)、安全率を \(n>1\) とすると \(\sigma_{\mathrm{allow}}=\sigma_{\mathrm{str}}/n\)。圧縮の座屈荷重はこの材料強度だけで決まらない。表の並びに従い本問の候補番号を選ぶ。

\Needspace{4\baselineskip}\paragraph{選択肢の検討（査読待ち初稿）}\mbox{}\par

\noindent\textbf{1：不適合} 引張・圧縮・剥離・安全率。横へたわむ圧縮不安定は座屈であり、層がはがれる剥離ではない。\par

\noindent\textbf{2：不適合} 圧縮・引張・剥離・安全率。薄い定規は引張より圧縮の安定性を失いやすく、AとBが逆である。横へたわむ圧縮不安定は座屈であり、層がはがれる剥離ではない。\par

\noindent\textbf{3：不適合} 圧縮・引張・剥離・ヤング率。薄い定規は引張より圧縮の安定性を失いやすく、AとBが逆である。横へたわむ圧縮不安定は座屈であり、層がはがれる剥離ではない。強度を割る無次元比は安全率であり、応力単位のヤング率ではない。\par

\noindent\textbf{4：適合（前提内）} 引張・圧縮・座屈・安全率。引張・圧縮・座屈・安全率が対応する。形の安定と材料強度は別の照査である。\par

\noindent\textbf{5：不適合} 圧縮・引張・座屈・ヤング率。薄い定規は引張より圧縮の安定性を失いやすく、AとBが逆である。強度を割る無次元比は安全率であり、応力単位のヤング率ではない。\par

\paragraph{再利用する考え方} 材料の強さ、構造の安定、変形による機能低下を分ける。寸法だけを変えるとき何を固定したかを明示し、同じ理論でも候補番号を年度間で流用しない。

\paragraph{共通理論} \hyperref[theory:structural:buckling]{第\ref*{theory:structural:buckling}節}, \hyperref[theory:structural:allowable]{第\ref*{theory:structural:allowable}節}

\par 参考：\cite{struct_mit_beam,struct_mit_buckling,struct_mlit_limit}。

\paragraph{関連問題} \hyperref[q:h25-reg-basic-1-2]{H25-REG-BASIC-1-2}

\par 難易度 2（校準初稿）、優先度未評価。

\subsection{令和元年度 問1-4}
\label{q:r01-reg-basic-1-4}

\noindent\textbf{原典：} sources/R01/basic.pdf、PDF物理4頁。

\textbf{公式正答：} 5（single）。

\paragraph{問題の要旨} 細長い棒の軸圧縮による横変形を、静的座屈として説明する。

\Needspace{4\baselineskip}\paragraph{確認した条件・記述（独自要約）}\mbox{}\par

\noindent\textbf{判断方向：} 空欄の適切な語・数値又は正誤の組合せ。全候補の順序は各肢に記録。\par

\noindent\textbf{原典照合：} 16原典画像で17問の全候補、図、数値、不等号、表の列順を確認。\par

\noindent\textbf{前提：} 解法に記録した弾性範囲、荷重条件、幾何、独立性を明示する。\par

\paragraph{問われる核心} 細長い棒の軸圧縮による横変形を、静的座屈として説明する。

\paragraph{必要知識} 圧縮荷重、曲げ剛性EI、端部拘束に依存する有効長KL、座屈と共振の差。

\paragraph{考え方・正答を選ぶ理由（査読待ち初稿）} 軸圧縮のもとで横たわみを増やす効果が曲げの復元に勝つと直線状態が不安定になる。\(P_{\mathrm{cr}}=\pi^2EI/(KL)^2\) であり、断面円の \(I=\pi d^4/64\) より、小径ほど臨界荷重が小さい。図の現象はねじれや振動の共振ではなく曲げ座屈である。理想棒では臨界を超える条件を選び、実物の初期不整に対しては臨界前からたわみが生じ得る点も分ける。

\Needspace{4\baselineskip}\paragraph{選択肢の検討（査読待ち初稿）}\mbox{}\par

\noindent\textbf{1：不適合} 小さい・以下・ねじれ・座屈。理想的な直線平衡の安定性を失うのは臨界荷重を超える側である。描かれた横変位は曲げであり、断面の軸回り回転であるねじれではない。\par

\noindent\textbf{2：不適合} 大きい・以下・ねじれ・共振。同じ長さ・材質では小径ほど曲げ剛性EIが小さく座屈しやすい。理想的な直線平衡の安定性を失うのは臨界荷重を超える側である。描かれた横変位は曲げであり、断面の軸回り回転であるねじれではない。静的圧縮での安定性喪失は座屈であり、周期的外力と固有振動の共振ではない。\par

\noindent\textbf{3：不適合} 小さい・超える・ねじれ・共振。描かれた横変位は曲げであり、断面の軸回り回転であるねじれではない。静的圧縮での安定性喪失は座屈であり、周期的外力と固有振動の共振ではない。\par

\noindent\textbf{4：不適合} 大きい・超える・曲げ・共振。同じ長さ・材質では小径ほど曲げ剛性EIが小さく座屈しやすい。静的圧縮での安定性喪失は座屈であり、周期的外力と固有振動の共振ではない。\par

\noindent\textbf{5：適合（前提内）} 小さい・超える・曲げ・座屈。小さい径・臨界荷重を超える・曲げ・座屈の対応である。\par

\paragraph{再利用する考え方} 材料の強さ、構造の安定、変形による機能低下を分ける。寸法だけを変えるとき何を固定したかを明示し、同じ理論でも候補番号を年度間で流用しない。

\paragraph{共通理論} \hyperref[theory:structural:buckling]{第\ref*{theory:structural:buckling}節}

\par 参考：\cite{struct_mit_beam,struct_mit_buckling,struct_mlit_limit}。

\paragraph{関連問題} \hyperref[q:r05-reg-basic-1-2]{R05-REG-BASIC-1-2}

\par 難易度 2（校準初稿）、優先度未評価。

\subsection{令和3年度 問1-5}
\label{q:r03-reg-basic-1-5}

\noindent\textbf{原典：} sources/R03/basic.pdf、PDF物理4頁。

\textbf{公式正答：} 5（single）。

\paragraph{問題の要旨} オイラー座屈の長さ・強度依存性と、許容応力・限界状態を組合せで検討する。

\Needspace{4\baselineskip}\paragraph{確認した条件・記述（独自要約）}\mbox{}\par

\noindent\textbf{判断方向：} 空欄の適切な語・数値又は正誤の組合せ。全候補の順序は各肢に記録。\par

\noindent\textbf{原典照合：} 16原典画像で17問の全候補、図、数値、不等号、表の列順を確認。\par

\noindent\textbf{前提：} 解法に記録した弾性範囲、荷重条件、幾何、独立性を明示する。\par

\paragraph{問われる核心} オイラー座屈の長さ・強度依存性と、許容応力・限界状態を組合せで検討する。

\paragraph{必要知識} オイラー式の弾性条件、Lの逆二乗、ヤング率と降伏強度の差、安全率の除算、複数の限界状態。

\paragraph{考え方・正答を選ぶ理由（査読待ち初稿）} \(P_{\mathrm{cr}}\propto L^{-2}\) なので長い柱ほど座屈荷重は小さい。弾性座屈域なら同じ \(E,I,KL\) の柱の臨界荷重は降伏強度を直接含まない。ただし算出座屈応力が降伏域に入る柱には弾性式をそのまま適用しない。許容応力は強度を安全率で割る。使用・終局・疲労等の限界状態を分けるため、誤正誤誤の5である。

\Needspace{4\baselineskip}\paragraph{選択肢の検討（査読待ち初稿）}\mbox{}\par

\noindent\textbf{1：不適合} 正誤正正。アは長いほど臨界荷重が小さくなるので誤。イを誤、ウ・エを正とする点も違う。\par

\noindent\textbf{2：不適合} 正正誤正。アは逆二乗則に反し、エは限界状態を一つだけに限定している。\par

\noindent\textbf{3：不適合} 誤誤誤正。イは弾性オイラー式の条件内では同じE・幾何・拘束で等しい座屈荷重となるため正。エも誤。\par

\noindent\textbf{4：不適合} 誤正正誤。ウの強度に安全率を掛けた許容応力は逆である。\par

\noindent\textbf{5：適合（前提内）} 誤正誤誤。アは長さ依存の方向が誤、イは弾性オイラー範囲で正、ウは除算でなく乗算のため誤、エは限界状態の単一化のため誤。\par

\paragraph{再利用する考え方} 材料の強さ、構造の安定、変形による機能低下を分ける。寸法だけを変えるとき何を固定したかを明示し、同じ理論でも候補番号を年度間で流用しない。

\paragraph{共通理論} \hyperref[theory:structural:buckling]{第\ref*{theory:structural:buckling}節}, \hyperref[theory:structural:allowable]{第\ref*{theory:structural:allowable}節}, \hyperref[theory:structural:limit]{第\ref*{theory:structural:limit}節}

\par 参考：\cite{struct_mit_beam,struct_mit_buckling,struct_mlit_limit}。

\paragraph{関連問題} \hyperref[q:h25-reg-basic-1-2]{H25-REG-BASIC-1-2}, \hyperref[q:r06-reg-basic-1-2]{R06-REG-BASIC-1-2}

\par 難易度 2（校準初稿）、優先度未評価。

\subsection{令和5年度 問1-2}
\label{q:r05-reg-basic-1-2}

\noindent\textbf{原典：} sources/R05/basic.pdf、PDF物理2頁。

\textbf{公式正答：} 4（single）。

\paragraph{問題の要旨} 細長い棒の軸圧縮による横変形を、静的座屈として説明する。

\Needspace{4\baselineskip}\paragraph{確認した条件・記述（独自要約）}\mbox{}\par

\noindent\textbf{判断方向：} 空欄の適切な語・数値又は正誤の組合せ。全候補の順序は各肢に記録。\par

\noindent\textbf{原典照合：} 16原典画像で17問の全候補、図、数値、不等号、表の列順を確認。\par

\noindent\textbf{前提：} 解法に記録した弾性範囲、荷重条件、幾何、独立性を明示する。\par

\paragraph{問われる核心} 細長い棒の軸圧縮による横変形を、静的座屈として説明する。

\paragraph{必要知識} 圧縮荷重、曲げ剛性EI、端部拘束に依存する有効長KL、座屈と共振の差。

\paragraph{考え方・正答を選ぶ理由（査読待ち初稿）} 軸圧縮のもとで横たわみを増やす効果が曲げの復元に勝つと直線状態が不安定になる。\(P_{\mathrm{cr}}=\pi^2EI/(KL)^2\) であり、断面円の \(I=\pi d^4/64\) より、小径ほど臨界荷重が小さい。図の現象はねじれや振動の共振ではなく曲げ座屈である。理想棒では臨界を超える条件を選び、実物の初期不整に対しては臨界前からたわみが生じ得る点も分ける。

\Needspace{4\baselineskip}\paragraph{選択肢の検討（査読待ち初稿）}\mbox{}\par

\noindent\textbf{1：不適合} 大きい・以下・ねじれ・共振。同じ長さ・材質では小径ほど曲げ剛性EIが小さく座屈しやすい。理想的な直線平衡の安定性を失うのは臨界荷重を超える側である。描かれた横変位は曲げであり、断面の軸回り回転であるねじれではない。静的圧縮での安定性喪失は座屈であり、周期的外力と固有振動の共振ではない。\par

\noindent\textbf{2：不適合} 小さい・超える・ねじれ・座屈。描かれた横変位は曲げであり、断面の軸回り回転であるねじれではない。\par

\noindent\textbf{3：不適合} 大きい・超える・曲げ・共振。同じ長さ・材質では小径ほど曲げ剛性EIが小さく座屈しやすい。静的圧縮での安定性喪失は座屈であり、周期的外力と固有振動の共振ではない。\par

\noindent\textbf{4：適合（前提内）} 小さい・超える・曲げ・座屈。小さい径・臨界荷重を超える・曲げ・座屈の対応である。\par

\noindent\textbf{5：不適合} 小さい・以下・曲げ・共振。理想的な直線平衡の安定性を失うのは臨界荷重を超える側である。静的圧縮での安定性喪失は座屈であり、周期的外力と固有振動の共振ではない。\par

\paragraph{再利用する考え方} 材料の強さ、構造の安定、変形による機能低下を分ける。寸法だけを変えるとき何を固定したかを明示し、同じ理論でも候補番号を年度間で流用しない。

\paragraph{共通理論} \hyperref[theory:structural:buckling]{第\ref*{theory:structural:buckling}節}

\par 参考：\cite{struct_mit_beam,struct_mit_buckling,struct_mlit_limit}。

\paragraph{関連問題} \hyperref[q:r01-reg-basic-1-4]{R01-REG-BASIC-1-4}

\par 難易度 2（校準初稿）、優先度未評価。

\section{安全率と費用モデルの最適化}

\subsection{令和元年度再試験 問1-5}
\label{q:r01-ret-basic-1-5}

\noindent\textbf{原典：} sources/R01\_retest/basic.pdf、PDF物理5頁。

\textbf{公式正答：} 2（single）。

\paragraph{問題の要旨} 安全率に依存する損害費用と対策費用の和を、定義域で最小化する。

\Needspace{4\baselineskip}\paragraph{確認した条件・記述（独自要約）}\mbox{}\par

\noindent\textbf{判断方向：} 空欄の適切な語・数値又は正誤の組合せ。全候補の順序は各肢に記録。\par

\noindent\textbf{原典照合：} 16原典画像で17問の全候補、図、数値、不等号、表の列順を確認。\par

\noindent\textbf{前提：} 解法に記録した弾性範囲、荷重条件、幾何、独立性を明示する。\par

\paragraph{問われる核心} 安全率に依存する損害費用と対策費用の和を、定義域で最小化する。

\paragraph{必要知識} 与えられた費用関数、微分による停留点、二階微分の正値、定義域と単位。

\paragraph{考え方・正答を選ぶ理由（査読待ち初稿）} 費用の単位を億円にそろえる。\(x\ge1\)、\(C(x)=90/(1+4x)+10x\) である。\(C^\prime=-360/(1+4x)^2+10=0\) より \(1+4x=6\)、\(x=5/4=1.25\) で候補2となる。\(C^{\prime\prime}=2880/(1+4x)^3>0\) なので許容域内の唯一の最小である。候補の費用は順に28、27.5、約27.857、28.75、30億円である。

\Needspace{4\baselineskip}\paragraph{選択肢の検討（査読待ち初稿）}\mbox{}\par

\noindent\textbf{1：不適合} 1.00。与式に代入した総費用は約28億円であり、最小の候補より大きい。\par

\noindent\textbf{2：適合（前提内）} 1.25。与式に代入した総費用は約27.5億円であり、連続域での停留点かつ最小である。\par

\noindent\textbf{3：不適合} 1.50。与式に代入した総費用は約27.857142857億円であり、最小の候補より大きい。\par

\noindent\textbf{4：不適合} 1.75。与式に代入した総費用は約28.75億円であり、最小の候補より大きい。\par

\noindent\textbf{5：不適合} 2.00。与式に代入した総費用は約30億円であり、最小の候補より大きい。\par

\paragraph{再利用する考え方} 材料の強さ、構造の安定、変形による機能低下を分ける。寸法だけを変えるとき何を固定したかを明示し、同じ理論でも候補番号を年度間で流用しない。

\paragraph{共通理論} \hyperref[theory:structural:optimization]{第\ref*{theory:structural:optimization}節}

\paragraph{関連問題} \hyperref[q:r04-reg-basic-1-4]{R04-REG-BASIC-1-4}

\par 難易度 2（校準初稿）、優先度未評価。

\subsection{令和4年度 問1-4}
\label{q:r04-reg-basic-1-4}

\noindent\textbf{原典：} sources/R04/basic.pdf、PDF物理3頁。

\textbf{公式正答：} 1（single）。

\paragraph{問題の要旨} 安全率に依存する損害費用と対策費用の和を、定義域で最小化する。

\Needspace{4\baselineskip}\paragraph{確認した条件・記述（独自要約）}\mbox{}\par

\noindent\textbf{判断方向：} 空欄の適切な語・数値又は正誤の組合せ。全候補の順序は各肢に記録。\par

\noindent\textbf{原典照合：} 16原典画像で17問の全候補、図、数値、不等号、表の列順を確認。\par

\noindent\textbf{前提：} 解法に記録した弾性範囲、荷重条件、幾何、独立性を明示する。\par

\paragraph{問われる核心} 安全率に依存する損害費用と対策費用の和を、定義域で最小化する。

\paragraph{必要知識} 与えられた費用関数、微分による停留点、二階微分の正値、定義域と単位。

\paragraph{考え方・正答を選ぶ理由（査読待ち初稿）} 費用の単位を億円にそろえる。\(x>1\)、\(C(x)=9/(1+x)+x\) である。\(C^\prime=-9/(1+x)^2+1=0\) より \(1+x=3\)、\(x=2\) で候補1となる。\(C^{\prime\prime}=18/(1+x)^3>0\) なので許容域で唯一の最小である。候補の費用は順に5、約5.071、5.25、5.5、5.8億円である。

\Needspace{4\baselineskip}\paragraph{選択肢の検討（査読待ち初稿）}\mbox{}\par

\noindent\textbf{1：適合（前提内）} 2.0。与式に代入した総費用は約5億円であり、連続域での停留点かつ最小である。\par

\noindent\textbf{2：不適合} 2.5。与式に代入した総費用は約5.071428571億円であり、最小の候補より大きい。\par

\noindent\textbf{3：不適合} 3.0。与式に代入した総費用は約5.25億円であり、最小の候補より大きい。\par

\noindent\textbf{4：不適合} 3.5。与式に代入した総費用は約5.5億円であり、最小の候補より大きい。\par

\noindent\textbf{5：不適合} 4.0。与式に代入した総費用は約5.8億円であり、最小の候補より大きい。\par

\paragraph{再利用する考え方} 材料の強さ、構造の安定、変形による機能低下を分ける。寸法だけを変えるとき何を固定したかを明示し、同じ理論でも候補番号を年度間で流用しない。

\paragraph{共通理論} \hyperref[theory:structural:optimization]{第\ref*{theory:structural:optimization}節}

\paragraph{関連問題} \hyperref[q:r01-ret-basic-1-5]{R01-RET-BASIC-1-5}

\par 難易度 2（校準初稿）、優先度未評価。

\section{崩壊・共振・圧壊と構造モデル}

\subsection{令和2年度 問1-3}
\label{q:r02-reg-basic-1-3}

\noindent\textbf{原典：} sources/R02/basic.pdf、PDF物理3頁。

\textbf{公式正答：} 3（single）。

\par\textbf{検討未確定：} 短柱の圧縮が破断強度を超えるというエの語は材料・圧縮強度・降伏の条件が曖昧。延性材料が一定の引張破断強度で圧縮破断するとは一般化しない。初回3/公式3を保持。

\paragraph{問題の要旨} 静的崩壊、共振、座屈、短柱の圧縮破損と構造解析のモデルを区別する。

\Needspace{4\baselineskip}\paragraph{確認した条件・記述（独自要約）}\mbox{}\par

\noindent\textbf{判断方向：} 空欄の適切な語・数値又は正誤の組合せ。全候補の順序は各肢に記録。\par

\noindent\textbf{原典照合：} 16原典画像で17問の全候補、図、数値、不等号、表の列順を確認。\par

\noindent\textbf{前提：} 解法に記録した弾性範囲、荷重条件、幾何、独立性を明示する。\par

\paragraph{問われる核心} 静的崩壊、共振、座屈、短柱の圧縮破損と構造解析のモデルを区別する。

\paragraph{必要知識} 塑性崩壊、外力周波数と固有振動数、細長比、圧縮材料破損、構造解析モデル。

\paragraph{考え方・正答を選ぶ理由（査読待ち初稿）} アの崩壊は塑性ヒンジ等による荷重支持機構の喪失であり弾性範囲だけで決めない。イの共振は外力の周波数が固有振動数へ近づくときに起き、応力が振動数へ達するという比較は次元も違う。ウの細長い柱は圧縮座屈する。エは短く太い柱で材料圧縮破損が支配し得るという趣旨、オは荷重・支持条件と部材を表す解析モデルである。従って誤誤正正正、3とする。ただしエの破断強度を全材料の圧縮破壊に同じように適用しない。

\Needspace{4\baselineskip}\paragraph{選択肢の検討（査読待ち初稿）}\mbox{}\par

\noindent\textbf{1：不適合} 正正正誤誤。アは塑性崩壊を弾性で決める点が誤り、イも応力と振動数を比べている。エ・オを誤りとする点も違う。\par

\noindent\textbf{2：不適合} 誤正正正誤。イの応力が固有振動数へ達する説明は単位も機構も違う。オの解析モデルの定義を否定できない。\par

\noindent\textbf{3：適合（前提内）} 誤誤正正正。アは誤、イは誤、ウは圧縮座屈で正、エは短柱の材料圧縮破損を意図した正、オは部材と力・拘束のモデルで正となる。\par

\noindent\textbf{4：不適合} 正誤誤正正。アは塑性機構の崩壊荷重であって弾性だけでは決まらず、ウの圧縮座屈は正である。\par

\noindent\textbf{5：不適合} 正正誤誤正。ア・イは誤、ウは正なので最初の三記号が一致しない。\par

\paragraph{再利用する考え方} 材料の強さ、構造の安定、変形による機能低下を分ける。寸法だけを変えるとき何を固定したかを明示し、同じ理論でも候補番号を年度間で流用しない。

\paragraph{共通理論} \hyperref[theory:structural:modes]{第\ref*{theory:structural:modes}節}, \hyperref[theory:structural:buckling]{第\ref*{theory:structural:buckling}節}

\par 参考：\cite{struct_mit_beam,struct_mit_buckling,struct_mlit_limit}。

\paragraph{関連問題} \hyperref[q:h24-reg-basic-1-2]{H24-REG-BASIC-1-2}

\par 難易度 2（校準初稿）、優先度未評価。

\section{疲労・クリープと局所応力集中}

\subsection{令和4年度 問1-1}
\label{q:r04-reg-basic-1-1}

\noindent\textbf{原典：} sources/R04/basic.pdf、PDF物理1頁。

\textbf{公式正答：} 1（single）。

\par\textbf{検討未確定：} 疲労の平均応力依存は振幅等の同条件比較を要する。急変部の局所von Mises値は荷重/形状/拘束に依存し、常に上がる全称の証明ではない。公式/初回1を維持し条件を限定。

\paragraph{問題の要旨} 疲労の平均応力、クリープ温度、曲げ剛性と応力集中を区別する。

\Needspace{4\baselineskip}\paragraph{確認した条件・記述（独自要約）}\mbox{}\par

\noindent\textbf{判断方向：} 空欄の適切な語・数値又は正誤の組合せ。全候補の順序は各肢に記録。\par

\noindent\textbf{原典照合：} 16原典画像で17問の全候補、図、数値、不等号、表の列順を確認。\par

\noindent\textbf{前提：} 解法に記録した弾性範囲、荷重条件、幾何、独立性を明示する。\par

\paragraph{問われる核心} 疲労の平均応力、クリープ温度、曲げ剛性と応力集中を区別する。

\paragraph{必要知識} 応力振幅/平均応力、通常の平均応力補正の条件、高温クリープ、EIと変位、局所応力と公称応力。

\paragraph{考え方・正答を選ぶ理由（査読待ち初稿）} Aは通常の疲労限度線図の引張平均応力側である。同じ応力振幅なら平均引張応力の増加は一般に疲労寿命へ不利となるが、材料・残留応力・塑性域へ無条件には拡張しない。Bは高温での時間依存変形、Cは同じ荷重と断面の \(\delta\propto E^{-1}\) より小さいヤング率、Dは断面急変部の局所応力の上昇である。従って1を選ぶ。von Mises応力の局所値は荷重と形状に依存し、急変させればどの点でも必ず高くなるという命題にはしない。

\Needspace{4\baselineskip}\paragraph{選択肢の検討（査読待ち初稿）}\mbox{}\par

\noindent\textbf{1：適合（前提内）} 引張・高い・小さい・高い。振幅等を固定した通常の金属疲労モデルでは引張平均応力が不利、高温でクリープ、小Eで大変位、急変部で局所応力が高い。\par

\noindent\textbf{2：不適合} 引張・高い・大きい・小さい。同じ荷重と幾何では変位はEの逆数であり、小Eが大変位である。応力集中は局所上昇を表す。\par

\noindent\textbf{3：不適合} 圧縮・高い・小さい・小さい。一般的平均応力モデルの不利側は引張であり、応力集中を局所低下とはしない。\par

\noindent\textbf{4：不適合} 圧縮・高強度・小さい・高い。Aは引張側、Bは高温である。高強度を温度条件の代わりにはできない。\par

\noindent\textbf{5：不適合} 引張・高強度・大きい・高い。クリープのBは高温であり、曲げ変位のCは小さいEである。\par

\paragraph{再利用する考え方} 材料の強さ、構造の安定、変形による機能低下を分ける。寸法だけを変えるとき何を固定したかを明示し、同じ理論でも候補番号を年度間で流用しない。

\paragraph{共通理論} \hyperref[theory:structural:fatigue]{第\ref*{theory:structural:fatigue}節}, \hyperref[theory:structural:creep]{第\ref*{theory:structural:creep}節}, \hyperref[theory:structural:beam]{第\ref*{theory:structural:beam}節}, \hyperref[theory:structural:concentration]{第\ref*{theory:structural:concentration}節}

\par 参考：\cite{struct_nims_data,struct_nasa_mean,struct_mit_beam}。

\par 難易度 2（校準初稿）、優先度未評価。

\section{使用・終局・疲労の限界状態}

\subsection{平成24年度 問1-2}
\label{q:h24-reg-basic-1-2}

\noindent\textbf{原典：} sources/H24/basic.pdf、PDF物理2頁。

\textbf{公式正答：} 5（single）。

\par\textbf{検討未確定：} 全ての構成部材の破壊まで至らなければ安全性が十分という前文は、機構形成や全体座屈を除外しておらず一般には保証できない。公式5と初回5を保持し、余力の条件を限定する。

\paragraph{問題の要旨} 不静定構造の余力、部分損傷、限界状態と非線形解析の関係を整理する。

\Needspace{4\baselineskip}\paragraph{確認した条件・記述（独自要約）}\mbox{}\par

\noindent\textbf{判断方向：} 空欄の適切な語・数値又は正誤の組合せ。全候補の順序は各肢に記録。\par

\noindent\textbf{原典照合：} 16原典画像で17問の全候補、図、数値、不等号、表の列順を確認。\par

\noindent\textbf{前提：} 解法に記録した弾性範囲、荷重条件、幾何、独立性を明示する。\par

\paragraph{問われる核心} 不静定構造の余力、部分損傷、限界状態と非線形解析の関係を整理する。

\paragraph{必要知識} 静定と不静定、荷重再配分、塑性ヒンジ、材料・幾何非線形、終局/使用/疲労の限界。

\paragraph{考え方・正答を選ぶ理由（査読待ち初稿）} Aを高い、Bを限界状態、Cを部分的、Dを非線形性とする5である。余剰拘束のある構造で一部が降伏しても荷重を他部へ再配分できる場合がある。しかし塑性機構、全体座屈、接合部破壊が先に生じ得るため、「全ての部材が破壊しなければ安全」という前文を普遍命題にしない。

\Needspace{4\baselineskip}\paragraph{選択肢の検討（査読待ち初稿）}\mbox{}\par

\noindent\textbf{1：不適合} 高い・限界状態・部分的・線形性。損傷後の応答を線形のままとするDが違う。\par

\noindent\textbf{2：不適合} 低い・限界状態・全体的・非線形性。Aは余剰拘束が多い高不静定、Cは部分損傷を許す説明なので全体的とする点が違う。\par

\noindent\textbf{3：不適合} 高い・弾性・部分的・線形性。弾性設計に限定すると部分損傷後の余力を捉えず、Dも非線形である。\par

\noindent\textbf{4：不適合} 低い・弾性・全体的・線形性。高不静定の余力を低不静定へ置換し、弾性・全体的・線形へ限定している。\par

\noindent\textbf{5：適合（前提内）} 高い・限界状態・部分的・非線形性。想定された延性と荷重再配分のモデルでは高不静定、限界状態、部分的損傷、非線形が対応する。ただし全構成部材が破壊しないだけでは全体安全の証明にならない。\par

\paragraph{再利用する考え方} 材料の強さ、構造の安定、変形による機能低下を分ける。寸法だけを変えるとき何を固定したかを明示し、同じ理論でも候補番号を年度間で流用しない。

\paragraph{共通理論} \hyperref[theory:structural:limit]{第\ref*{theory:structural:limit}節}, \hyperref[theory:structural:redundancy]{第\ref*{theory:structural:redundancy}節}

\par 参考：\cite{struct_mit_beam,struct_mit_buckling,struct_mlit_limit}。

\paragraph{関連問題} \hyperref[q:r02-reg-basic-1-3]{R02-REG-BASIC-1-3}

\par 難易度 2（校準初稿）、優先度未評価。

\subsection{令和6年度 問1-2}
\label{q:r06-reg-basic-1-2}

\noindent\textbf{原典：} sources/R06/basic.pdf、PDF物理2頁。

\textbf{公式正答：} 2（single）。

\paragraph{問題の要旨} 通常使用、繰返し荷重、最大耐荷性能の三つの要求を限界状態へ対応させる。

\Needspace{4\baselineskip}\paragraph{確認した条件・記述（独自要約）}\mbox{}\par

\noindent\textbf{判断方向：} 空欄の適切な語・数値又は正誤の組合せ。全候補の順序は各肢に記録。\par

\noindent\textbf{原典照合：} 16原典画像で17問の全候補、図、数値、不等号、表の列順を確認。\par

\noindent\textbf{前提：} 解法に記録した弾性範囲、荷重条件、幾何、独立性を明示する。\par

\paragraph{問われる核心} 通常使用、繰返し荷重、最大耐荷性能の三つの要求を限界状態へ対応させる。

\paragraph{必要知識} 使用性、疲労による累積損傷、終局耐荷性能、規準別の分類粒度。

\paragraph{考え方・正答を選ぶ理由（査読待ち初稿）} 過大なたわみや振動等で通常の使用に支障が生じる状態は使用限界である。繰返し荷重の損傷が疲労限界、崩壊や不安定等で最大耐荷性能を失う状態が終局限界である。使用・疲労・終局の2を選ぶ。基準によっては疲労を終局に含める分類もあるが、本問の三項目対応を変えない。

\Needspace{4\baselineskip}\paragraph{選択肢の検討（査読待ち初稿）}\mbox{}\par

\noindent\textbf{1：不適合} 使用・終局・疲労。Bの繰返し荷重は疲労、Cの最大耐荷は終局であり入れ替わっている。\par

\noindent\textbf{2：適合（前提内）} 使用・疲労・終局。Aは通常の機能・快適性等の使用限界、Bは疲労、Cは終局であり対応する。\par

\noindent\textbf{3：不適合} 終局・疲労・使用。Aは使用、Cは終局なので逆転している。\par

\noindent\textbf{4：不適合} 疲労・使用・終局。Aは使用、Bは繰返しの疲労なので逆転している。\par

\noindent\textbf{5：不適合} 疲労・終局・使用。三項目とも用途との対応が一致しない。\par

\paragraph{再利用する考え方} 材料の強さ、構造の安定、変形による機能低下を分ける。寸法だけを変えるとき何を固定したかを明示し、同じ理論でも候補番号を年度間で流用しない。

\paragraph{共通理論} \hyperref[theory:structural:limit]{第\ref*{theory:structural:limit}節}

\par 参考：\cite{struct_mlit_limit}。

\paragraph{関連問題} \hyperref[q:r03-reg-basic-1-5]{R03-REG-BASIC-1-5}

\par 難易度 2（校準初稿）、優先度未評価。

\section{強度・荷重の同時分布と破壊確率}

\subsection{平成29年度 問1-6}
\label{q:h29-reg-basic-1-6}

\noindent\textbf{原典：} sources/H29/basic.pdf、PDF物理5頁。

\textbf{公式正答：} 2（single）。

\par\textbf{検討未確定：} 強度と荷重の独立性が原文に明記されない。周辺一様分布だけでは破壊確率を一意に決定できない。独立なら初回2/公式2、同じ周辺を持つ別結合で0又は2/15となる反例を保存する。

\paragraph{問題の要旨} 強度と荷重が一様分布するモデルで、荷重が強度を超える確率と独立性の役割を検討する。

\Needspace{4\baselineskip}\paragraph{確認した条件・記述（独自要約）}\mbox{}\par

\noindent\textbf{判断方向：} 空欄の適切な語・数値又は正誤の組合せ。全候補の順序は各肢に記録。\par

\noindent\textbf{原典照合：} 16原典画像で17問の全候補、図、数値、不等号、表の列順を確認。\par

\noindent\textbf{前提：} 解法に記録した弾性範囲、荷重条件、幾何、独立性を明示する。\par

\paragraph{問われる核心} 強度と荷重が一様分布するモデルで、荷重が強度を超える確率と独立性の役割を検討する。

\paragraph{必要知識} 同時分布と周辺分布、独立性、破壊事象S>R、密度の面積積分。

\paragraph{考え方・正答を選ぶ理由（査読待ち初稿）} 強度を \(R\sim U[18,23]\)、荷重を \(S\sim U[10,20]\) とする。独立なら \(f_{R,S}=1/(5\cdot10)\)。破壊条件 \(S>R\) の領域は \(18\le R\le20,\ R<S\le20\) であるから、\(P(S>R)=\int_{18}^{20}(20-r)/50\,dr=4/100=0.04\)、候補2である。ただし原文には独立の明記がない。共通の \(U\sim U[0,1]\) を用い、\(R=18+5U,S=10+10U\) なら破壊確率は0、\(R=18+5U,S=20-10U\) なら \(U<2/15\) より確率2/15である。いずれも同じ周辺一様分布を持つ。従って公式0.04は独立モデルを補った条件付きの解として示す。

\Needspace{4\baselineskip}\paragraph{選択肢の検討（査読待ち初稿）}\mbox{}\par

\noindent\textbf{1：不適合} 0.02。独立モデルの三角形面積をさらに半分にすると0.02となるが、その二重の半減は不要である。\par

\noindent\textbf{2：適合（前提内）} 0.04。独立を仮定すれば支持矩形の面積50に対し破壊三角形面積2なので0.04となる。\par

\noindent\textbf{3：不適合} 0.08。一辺2の正方形を全て破壊領域にすると0.08となる。対角線の半分だけが該当する。\par

\noindent\textbf{4：不適合} 0.10。0.10は独立モデルの領域積分と一致しない。周辺分布だけで確率が決まる値でもない。\par

\noindent\textbf{5：不適合} 0.20。強度の区間長5の逆数0.20は密度であり、破壊事象の確率ではない。\par

\paragraph{再利用する考え方} 材料の強さ、構造の安定、変形による機能低下を分ける。寸法だけを変えるとき何を固定したかを明示し、同じ理論でも候補番号を年度間で流用しない。

\paragraph{共通理論} \hyperref[theory:structural:reliability]{第\ref*{theory:structural:reliability}節}

\par 難易度 4（校準初稿）、優先度未評価。

\section{比強度と材料価格当たり性能}

\subsection{令和5年度 問1-1}
\label{q:r05-reg-basic-1-1}

\noindent\textbf{原典：} sources/R05/basic.pdf、PDF物理1頁。

\textbf{公式正答：} 5（single）。

\paragraph{問題の要旨} 与えられた鋼とCFRPの密度・強度・価格から、比強度とその価格当たり性能を比較する。

\Needspace{4\baselineskip}\paragraph{確認した条件・記述（独自要約）}\mbox{}\par

\noindent\textbf{判断方向：} 空欄の適切な語・数値又は正誤の組合せ。全候補の順序は各肢に記録。\par

\noindent\textbf{原典照合：} 16原典画像で17問の全候補、図、数値、不等号、表の列順を確認。\par

\noindent\textbf{前提：} 解法に記録した弾性範囲、荷重条件、幾何、独立性を明示する。\par

\paragraph{問われる核心} 与えられた鋼とCFRPの密度・強度・価格から、比強度とその価格当たり性能を比較する。

\paragraph{必要知識} 比強度σ/ρ、同じ質量の荷重支持能力、価格/比強度、所与価格と方向条件。

\paragraph{考え方・正答を選ぶ理由（査読待ち初稿）} 比強度は \(\sigma/\rho\) である。与えられた数値だけを使うと、CFRP/鋼の比強度比は \((2000/1600)/(400/7900)=24.6875\)、kg当たり価格比は \(16000/60=266.6667\)。単位比強度当たり価格比は \((16000/(2000/1600))/(60/(400/7900))\simeq10.8017\) で、最も近い約10、5を選ぶ。比強度が大きいことと費用当たり有利なことを区別する。価格は設問所与であり現在の相場ではない。CFRPの方向依存や圧縮特性を所与の引張強度へ無断で補わない。

\Needspace{4\baselineskip}\paragraph{選択肢の検討（査読待ち初稿）}\mbox{}\par

\noindent\textbf{1：不適合} 強度/密度・CFRP・2。比強度の定義と優れる材料は合うが、単位比強度当たり価格の倍率は約10.8であり2ではない。\par

\noindent\textbf{2：不適合} 密度/強度・CFRP・10。密度/強度は比強度の逆数である。\par

\noindent\textbf{3：不適合} 密度/強度・鋼・2。定義を逆転し、与条件での優れる材料と価格倍率も違う。\par

\noindent\textbf{4：不適合} 強度/密度・鋼・2。定義は合うが、強度/密度はCFRPの方が大きく、価格倍率も2ではない。\par

\noindent\textbf{5：適合（前提内）} 強度/密度・CFRP・10。強度/密度でCFRPが優れ、価格当たりの比較ではCFRPの単位比強度当たり費用が約10倍となる。\par

\paragraph{再利用する考え方} 材料の強さ、構造の安定、変形による機能低下を分ける。寸法だけを変えるとき何を固定したかを明示し、同じ理論でも候補番号を年度間で流用しない。

\paragraph{共通理論} \hyperref[theory:structural:specific]{第\ref*{theory:structural:specific}節}

\par 難易度 3（校準初稿）、優先度未評価。

\section{引張試験と公称応力・弾塑性}

\subsection{平成29年度 問1-4}
\label{q:h29-reg-basic-1-4}

\noindent\textbf{原典：} sources/H29/basic.pdf、PDF物理3頁。

\textbf{公式正答：} 5（single）。

\par\textbf{検討未確定：} 前文は荷重を断面積で割った値を応力とするが初期断面積を明示しない。最大値を引張強さとするには公称応力を用いる条件が必要。真応力の最大と同一視せず初回/公式5を保持。

\paragraph{問題の要旨} 引張試験の測定値と、弾性・塑性・引張強さの定義を対応させる。

\Needspace{4\baselineskip}\paragraph{確認した条件・記述（独自要約）}\mbox{}\par

\noindent\textbf{判断方向：} 空欄の適切な語・数値又は正誤の組合せ。全候補の順序は各肢に記録。\par

\noindent\textbf{原典照合：} 16原典画像で17問の全候補、図、数値、不等号、表の列順を確認。\par

\noindent\textbf{前提：} 解法に記録した弾性範囲、荷重条件、幾何、独立性を明示する。\par

\paragraph{問われる核心} 引張試験の測定値と、弾性・塑性・引張強さの定義を対応させる。

\paragraph{必要知識} 初期断面積と初期標点距離、荷重/伸び、弾性/塑性、延性と公称引張強さ。

\paragraph{考え方・正答を選ぶ理由（査読待ち初稿）} 測定する荷重 \(F\) と伸び \(\Delta L\) から、公称応力 \(F/A_0\) と公称ひずみ \(\Delta L/L_0\) を求める。除荷で回復する変形が弾性、残る変形が塑性である。単調引張試験での \(\max F/A_0\) が引張強さであり、破断時の応力、繰返し荷重での疲労限度、真応力 \(F/A(t)\) と分ける。従って伸び・弾性・塑性・引張強さを選ぶ。

\Needspace{4\baselineskip}\paragraph{選択肢の検討（査読待ち初稿）}\mbox{}\par

\noindent\textbf{1：不適合} ひずみ・弾性・延性・疲労限度。直接測る長さの増分は伸びであり、ひずみは初期長さで除した値である。残る変形の性質は塑性であり、延性は破断まで塑性変形を許す能力である。単調引張の最大公称応力は引張強さであり、疲労限度や破断時応力と区別する。\par

\noindent\textbf{2：不適合} 伸び・塑性・弾性・引張強さ。除荷後に回復する性質は弾性である。残る変形の性質は塑性であり、延性は破断まで塑性変形を許す能力である。\par

\noindent\textbf{3：不適合} 伸び・弾性・延性・疲労限度。残る変形の性質は塑性であり、延性は破断まで塑性変形を許す能力である。単調引張の最大公称応力は引張強さであり、疲労限度や破断時応力と区別する。\par

\noindent\textbf{4：不適合} ひずみ・延性・塑性・破断強度。直接測る長さの増分は伸びであり、ひずみは初期長さで除した値である。除荷後に回復する性質は弾性である。単調引張の最大公称応力は引張強さであり、疲労限度や破断時応力と区別する。\par

\noindent\textbf{5：適合（前提内）} 伸び・弾性・塑性・引張強さ。伸び・弾性・塑性・引張強さが対応する。公称応力の断面積は初期値である。\par

\paragraph{再利用する考え方} 材料の強さ、構造の安定、変形による機能低下を分ける。寸法だけを変えるとき何を固定したかを明示し、同じ理論でも候補番号を年度間で流用しない。

\paragraph{共通理論} \hyperref[theory:structural:tensile]{第\ref*{theory:structural:tensile}節}

\par 参考：\cite{struct_mit_beam,struct_mit_buckling,struct_mlit_limit}。

\paragraph{関連問題} \hyperref[q:r05-reg-basic-1-3]{R05-REG-BASIC-1-3}

\par 難易度 2（校準初稿）、優先度未評価。

\subsection{令和5年度 問1-3}
\label{q:r05-reg-basic-1-3}

\noindent\textbf{原典：} sources/R05/basic.pdf、PDF物理3頁。

\textbf{公式正答：} 3（single）。

\paragraph{問題の要旨} 引張試験の測定値と、弾性・塑性・引張強さの定義を対応させる。

\Needspace{4\baselineskip}\paragraph{確認した条件・記述（独自要約）}\mbox{}\par

\noindent\textbf{判断方向：} 空欄の適切な語・数値又は正誤の組合せ。全候補の順序は各肢に記録。\par

\noindent\textbf{原典照合：} 16原典画像で17問の全候補、図、数値、不等号、表の列順を確認。\par

\noindent\textbf{前提：} 解法に記録した弾性範囲、荷重条件、幾何、独立性を明示する。\par

\paragraph{問われる核心} 引張試験の測定値と、弾性・塑性・引張強さの定義を対応させる。

\paragraph{必要知識} 初期断面積と初期標点距離、荷重/伸び、弾性/塑性、延性と公称引張強さ。

\paragraph{考え方・正答を選ぶ理由（査読待ち初稿）} 測定する荷重 \(F\) と伸び \(\Delta L\) から、公称応力 \(F/A_0\) と公称ひずみ \(\Delta L/L_0\) を求める。除荷で回復する変形が弾性、残る変形が塑性である。単調引張試験での \(\max F/A_0\) が引張強さであり、破断時の応力、繰返し荷重での疲労限度、真応力 \(F/A(t)\) と分ける。従って伸び・弾性・塑性・引張強さを選ぶ。

\Needspace{4\baselineskip}\paragraph{選択肢の検討（査読待ち初稿）}\mbox{}\par

\noindent\textbf{1：不適合} ひずみ・弾性・延性・疲労限度。直接測る長さの増分は伸びであり、ひずみは初期長さで除した値である。残る変形の性質は塑性であり、延性は破断まで塑性変形を許す能力である。単調引張の最大公称応力は引張強さであり、疲労限度や破断時応力と区別する。\par

\noindent\textbf{2：不適合} 伸び・塑性・弾性・引張強さ。除荷後に回復する性質は弾性である。残る変形の性質は塑性であり、延性は破断まで塑性変形を許す能力である。\par

\noindent\textbf{3：適合（前提内）} 伸び・弾性・塑性・引張強さ。伸び・弾性・塑性・引張強さが対応する。公称応力の断面積は初期値である。\par

\noindent\textbf{4：不適合} 伸び・弾性・延性・疲労限度。残る変形の性質は塑性であり、延性は破断まで塑性変形を許す能力である。単調引張の最大公称応力は引張強さであり、疲労限度や破断時応力と区別する。\par

\noindent\textbf{5：不適合} ひずみ・延性・塑性・引張強さ。直接測る長さの増分は伸びであり、ひずみは初期長さで除した値である。除荷後に回復する性質は弾性である。\par

\paragraph{再利用する考え方} 材料の強さ、構造の安定、変形による機能低下を分ける。寸法だけを変えるとき何を固定したかを明示し、同じ理論でも候補番号を年度間で流用しない。

\paragraph{共通理論} \hyperref[theory:structural:tensile]{第\ref*{theory:structural:tensile}節}

\par 参考：\cite{struct_mit_beam,struct_mit_buckling,struct_mlit_limit}。

\paragraph{関連問題} \hyperref[q:h29-reg-basic-1-4]{H29-REG-BASIC-1-4}

\par 難易度 2（校準初稿）、優先度未評価。

\section{ユニバーサルデザイン}

\subsection{平成24年度 問1-4}
\label{q:h24-reg-basic-1-4}

\noindent\textbf{原典：} sources/H24/basic.pdf、PDF物理4頁。

\textbf{公式正答：} 5（single）。

\par\textit{原典本文の校正と独自解説は未作成。分類のみの登録項目。}

\subsection{平成26年度 問1-1}
\label{q:h26-reg-basic-1-1}

\noindent\textbf{原典：} sources/H26/basic.pdf、PDF物理1頁。

\textbf{公式正答：} 5（single）。

\par\textit{原典本文の校正と独自解説は未作成。分類のみの登録項目。}

\subsection{平成30年度 問1-3}
\label{q:h30-reg-basic-1-3}

\noindent\textbf{原典：} sources/H30/basic.pdf、PDF物理3頁。

\textbf{公式正答：} 2（single）。

\par\textit{原典本文の校正と独自解説は未作成。分類のみの登録項目。}

\subsection{令和2年度 問1-1}
\label{q:r02-reg-basic-1-1}

\noindent\textbf{原典：} sources/R02/basic.pdf、PDF物理1頁。

\textbf{公式正答：} 3（single）。

\par\textit{原典本文の校正と独自解説は未作成。分類のみの登録項目。}

\subsection{令和3年度 問1-1}
\label{q:r03-reg-basic-1-1}

\noindent\textbf{原典：} sources/R03/basic.pdf、PDF物理1頁。

\textbf{公式正答：} 4（single）。

\par\textit{原典本文の校正と独自解説は未作成。分類のみの登録項目。}

\subsection{令和6年度 問1-1}
\label{q:r06-reg-basic-1-1}

\noindent\textbf{原典：} sources/R06/basic.pdf、PDF物理1頁。

\textbf{公式正答：} 1（single）。

\par\textit{原典本文の校正と独自解説は未作成。分類のみの登録項目。}

\section{未確定}

\subsection{令和元年度 問1-2}
\label{q:r01-reg-basic-1-2}

\noindent\textbf{原典：} sources/R01/basic.pdf、PDF物理2頁。

\textbf{公式正答：} 2（single）。

\par\textit{原典本文の校正と独自解説は未作成。分類のみの登録項目。}

\subsection{令和3年度 問1-6}
\label{q:r03-reg-basic-1-6}

\noindent\textbf{原典：} sources/R03/basic.pdf、PDF物理5頁。

\textbf{公式正答：} 3（single）。

\par\textit{原典本文の校正と独自解説は未作成。分類のみの登録項目。}

\subsection{令和4年度 問1-2}
\label{q:r04-reg-basic-1-2}

\noindent\textbf{原典：} sources/R04/basic.pdf、PDF物理2頁。

\textbf{公式正答：} 4（single）。

\par\textit{原典本文の校正と独自解説は未作成。分類のみの登録項目。}

\subsection{令和4年度 問1-6}
\label{q:r04-reg-basic-1-6}

\noindent\textbf{原典：} sources/R04/basic.pdf、PDF物理4頁。

\textbf{公式正答：} 2（single）。

\par\textit{原典本文の校正と独自解説は未作成。分類のみの登録項目。}

\subsection{令和6年度 問1-4}
\label{q:r06-reg-basic-1-4}

\noindent\textbf{原典：} sources/R06/basic.pdf、PDF物理3頁。

\textbf{公式正答：} 3（single）。

\par\textit{原典本文の校正と独自解説は未作成。分類のみの登録項目。}

\subsection{令和6年度 問1-6}
\label{q:r06-reg-basic-1-6}

\noindent\textbf{原典：} sources/R06/basic.pdf、PDF物理5頁。

\textbf{公式正答：} 3（single）。

\par\textit{原典本文の校正と独自解説は未作成。分類のみの登録項目。}
